
%Updated and Revised Monthly Template 5.7.2024 by Bonnie Ponce and Bev Ruedi copyright the Mathematical Association of America and all rights reserved.

\documentclass{amsart}
%\usepackage{maa-monthly}
\usepackage{mathdots}
\usepackage{bbm}
\usepackage{epsfig}
\usepackage{epstopdf}


\newtheorem{theorem}{Theorem}
\newtheorem{proposition}[theorem]{Proposition}
\newtheorem{lemma}[theorem]{Lemma}
\newtheorem{corollary}[theorem]{Corollary}

\theoremstyle{definition}
\newtheorem{definition}[theorem]{Definition}
\newtheorem{remark}[theorem]{Remark}
\newtheorem{conjecture}[theorem]{Conjecture}

\usepackage{srcltx,amssymb}
\usepackage{graphicx}
\usepackage{amsthm}
\raggedbottom
\usepackage{amsmath}
\usepackage{amscd}
\usepackage{slashed}
\numberwithin{theorem}{section} 
\numberwithin{equation}{section} 
  
\begin{document} 

    
    \title{Anti-rotors and algebra invariants}
    
       \author{Fred Greensite}
 
       \begin{abstract} We derive previously unrecognized geometrically-based isomorphism invariants for real finite-dimensional unital associative algebras, resulting from examination of transformations from an algebra to itself that are dependent on the algebra's bilinear product operation.  
   \end{abstract}
   
    \maketitle
    
    \markboth{F. Greensite }{Anti-rotors and algebra invariants}

\thispagestyle{empty}
\section{Introduction}
As the means of understanding algebra structure, identification of the invariants of an algebra is a task of fundamental importance as well as an open-ended endeavor given that the algebra classification problem is apparently intractable \cite{eom:2024,belitskii:2005}.  While classical algebra isomorphism invariants include things like an algebra's dimension, types of ideals, and invariants stemming from the Jordan-H\"older Theorem, more recently they have derived from module categories, e.g., involving homological formalisms such as found in \cite{hochschild:1945,gerstenhaber:1964,loday:1998}.  The invariants we present here belong to a different class.  

\bigskip

{\it Conventions.} Unless stated otherwise, use of the word ``algebra" always refers to a unital associative algebra with vector space of elements $\mathbb{R}^n$, $n\ge 1$.  The standard topology on $\mathbb{R}^n$ is assumed, as well as the standard basis with elements taken to be column vectors.  Superscript $^T$ indicates dual vector or matrix transpose as applicable, and ``$\,\cdot\,$" is the usual dot product.  


\section{The Skewer and the anti-rotor} Our program begins with consideration of the transformations of an algebra $A$ to itself that are dependent on its bilinear product operation, as a means of studying the algebra.   Each such unary operation, $f:D_{A,f}\rightarrow A$, is a function as well as vector field (e.g., $f(s)=s^{-1}$, $f(s)=s^2$, etc.), where $D_{A,f}$ is an appropriate domain in the algebra's vector space of elements.  In examination of a vector field it is of interest to consider the associated infinitesimal rotation tensor (the exterior derivative of the differential 1-form dual to the vector field; its Hodge dual is the curl of the vector field).  Its computation requires choice of a metric (i.e., deriving from a symmetric bilinear form, and also referred to as a scalar product or metric tensor).  Specification of the standard topology entails imposition of the Euclidean inner product as the metric, allowing the computation to proceed.  However, this choice of metric is arbitrary (i.e., not inherent in the axioms of an algebra), and thus can only be considered a default against which the results of all other choices should be compared.  Hence, with respect to a specific choice of function $f$ to study, we are motivated to examine the mapping, \enlargethispage{1cm}  \begin{equation}\label{1/23/25.1} \slashed{\text{S}}_f: \text{Sym}_n \rightarrow \delta\text{Rot}_n,\end{equation}  
with, \begin{equation}\label{1/23/25.2} \slashed{\text{S}}_f[\mu] = d\left((f(s))^T \mu \,\mathbf{d}s\right).  \end{equation}
In (\ref{1/23/25.1}), $\text{Sym}_n$ is the linear space of all possible metric tensors relevant to $\mathbb{R}^n$ (inclusive of metrics that are not necessarily positive-definite and not necessarily nondegenerate), and $\delta\text{Rot}_n$ is the linear space of all possible infinitesimal rotation tensors relevant to $\mathbb{R}^n$.  Equation (\ref{1/23/25.2}) is interpreted as follows.  An algebra element is understood to be a column vector $s=(x_1,\dots,x_n)^T$, we define $\mathbf{d}s \equiv (dx_1,\dots,dx_n)^T$, and $\mu\in \text{Sym}_n$ is thereby a real symmetric $(n\times n)$-matrix.  On the right-hand-side of (\ref{1/23/25.2}), $d$ represents the exterior derivative operation, thus acting on the 1-form $(f(s))^T \mu \,\mathbf{d}s = \big(\mu f(s)\big)\cdot\mathbf{d}s$.   Consequently, the right-hand-side of (\ref{1/23/25.2}) yields the infinitesimal rotation tensor associated with the vector field $f(s)$ with respect to the metric $\mu$ - or, equivalently, the infinitesimal rotation tensor of vector field $\mu f(s)$ with respect to the Euclidean metric.


$\slashed{\text{S}}_f$ is a linear transformation mapping symmetric tensors to skew-symmetric tensors, so we refer to $\slashed{\text{S}}_f$ as a ``Skewer" (the slash through ``symmetric" S represents a skewer, as with a shish-kabob). Since we assume the standard basis, $\text{Sym}_n$ can be represented by a parametrized symmetric $(n\times n)$-matrix $M$ with $\frac{n(n+1)}{2}$ independent parameters, i.e., it is given as a symmetric matrix with a different real parameter in each of its upper triangular entries.  Every possible metric tensor is given by some realization of the set of parameters (a realization being the result of any particular choice of $\frac{n(n+1)}{2}$ real values for the members of the set of parameters).  Consequently, $M$ is mapped by a Skewer to a parametrized skew-symmetric matrix representing the resulting infinitesimal rotation tensors.  Each entry not on the main diagonal of the latter matrix is a rational function of the $n$ components of the elements of algebra $A$, and the above $\frac{n(n+1)}{2}$ parameters appear as coefficients of constituent rational function terms which sum to the entry.  

While the image of a Skewer is thereby complicated, the kernel of a Skewer is much less so for one and possibly two reasons.  First, the components of the elements of $A$ do not appear.  Second, there is a lot of computational evidence that kernel of the Skewer is independent of the choice of $f$ over the class of functions that are dependent on the algebra's product operation.   Note that for $f$ {\it not} dependent on the algebra product (i.e., of the form $f(s) = \alpha s^0 + \beta s^1$, $\alpha,\beta\in\mathbb{R}$), the kernel of the Skewer is simply $\text{Sym}_n$ itself.  

We will henceforth adopt ``anti-rotor" as a more descriptive term for the kernel of the Skewer, because each of the latter's members annihilates the curl of the relevant vector fields. 
For definiteness, in the first seven sections of this paper we will work with $f(s)=s^{-1}$, since this choice is associated with a productive ``dual" of the anti-rotor.  The issue of the anti-rotor's independence from the algebra-dependent vector fields is considered in Section 8.

Alternative motivations for our program, from Euclidean geometry and Inverse Problems respectively, are found in \cite{fgreensite:2025b} and \cite{fgreensite:2025c}.


\enlargethispage{1.5cm}
\section{The anti-rotor and resulting invariants} The following will serve until Section 8. 
     \begin{definition}\label{7/21/23.1} With respect to elements $s$ of a given $n$-dimensional algebra $A$ written in the standard basis as $s= (x_1,\dots,x_n)^T$, where these elements are in an open ball in $\mathbb{R}^n$ consisting only of units and whose central point is the multiplicative identity element $\mathbf{1}_A$, an {\it uncurling metric $L$} is a symmetric member of the full algebra of real $(n\times n)$-matrices $M_n(\mathbb{R})$,  satisfying \begin{equation}\label{5/22/25.3} d\left((Ls^{-1})\cdot \mathbf{d}s\right) = d\left((s^{-1})^T L \mathbf{d}s\right) = 0, \end{equation} where $\mathbf{d}s \equiv (dx_1,\dots,dx_n)^T$.  The real vector space of uncurling metrics is the algebra's {\it anti-rotor}, $\mathfrak{u}_A$.  It may also be denoted by $\mathfrak{u}$ if no confusion arises.    \end{definition}
     
Using the left regular representation, the existence of an open ball centered at the algebra's multiplicative identity that is devoid of non-units is easily established. 

Many things stem from the following.

  \begin{theorem} \label{6/6/23.1} If $K:A_1\rightarrow A_2$ is an algebra isomorphism, then \begin{equation}\label{5/23/23.1} \mathfrak{u}_{A_1} = K^T \mathfrak{u}_{A_2} K, \end{equation} 
   where $K^T \mathfrak{u}_{A_2} K \equiv \{K^T L K:L\in \mathfrak{u}_{A_2}\}$.
   \end{theorem}

\begin{proof} By hypothesis, a nonsingular linear change of variable expression for the elements of $A_2$ in terms of the elements of $A_1$ can be written as $s_2= Ks_1$. \mbox{For $L_2\in \mathfrak{u}_{A_2}$} we obtain, \begin{equation} 0 = d\left((s_2^{-1})^T L_2 \mathbf{d}s_2\right) =  d\left(((Ks_1)^{-1})^T L_2 \mathbf{d}(Ks_1)\right) = d\left((s_1^{-1})^T \big[K^T L_2K\big]\mathbf{d}s_1\right), \label{9/21/23.1}  \end{equation} where we have used Definition \ref{7/21/23.1}, the commutativity of a differential with a linear transformation, and $(Ks_1)^{-1} = K(s_1^{-1})$ (since $K$ is an isomorphism).   It follows from (\ref{9/21/23.1}) that every uncurling metric  of algebra $A_2$ is associated with an uncurling metric of $A_1$, i.e., from the last equality in (\ref{9/21/23.1}) we have  $K^TL_2K\in\mathfrak{u}_{A_1}$. The implied map from $\mathfrak{u}_{A_2}$ to $\mathfrak{u}_{A_1}$ is obviously injective since $K^TL_aK=K^TL_bK$ implies $L_a=L_b$ due to the invertibility of $K$.   The isomorphism hypothesis allows us to apply an analogous argument employing $K^{-1}$, indicating that for $L_1\in\mathfrak{u}_{A_1}$ we have the associated $(K^{-1})^T L_1 K^{-1} \in \mathfrak{u}_{A_2}$, implying an injective map from $\mathfrak{u}_{A_1}$ to $\mathfrak{u}_{A_2}$.  The theorem then follows from the bijection implied by the Cantor-Schoeder-Bernstein Theorem. 
\end{proof}

\begin{corollary}\label{7/31/25.1} If $A_1$ and $A_2$ are isomorphic algebras, then $\mathfrak{u}_{A_1}$ and $\mathfrak{u}_{A_2}$ are the same up to a change \mbox{of basis}. \end{corollary}  

Although it will be shown in Section 6 that anti-rotor dimension is generally greater than 1, for now we are content with,       

\begin{proposition}\label{6/11/25.1} The dimension of an algebra's anti-rotor is at least 1. \end{proposition} 
\begin{proof} For invertible elements $s = (s_{ij})\in M_n(\mathbb{R})$, with det$(s)$ being the determinant of $s$, we have the well-known,
  \begin{equation}\label{11/5/21.1}[\nabla \text{det}(s)] = \big[{\rm Adj}(s)\big]^T = \mbox{det}(s)[s^{-1}]^T,\end{equation} where Adj$(s)$ is the adjugate matrix, and the left-hand-side is considered to be a matrix according to $[\nabla \text{det}(s)]\equiv \left(\frac{\partial\,\text{det}(s)}{\partial s_{ij}}\right)$. The first of these equalities is a special case of Jacobi's formula, and the second follows from the classical matrix inverse expression resulting from Laplace expansion. 
    But equivalently, (\ref{11/5/21.1}) can be understood as equating three vectors, where the matrix format of each bracketed expression is identified with the vector given by the ordered \mbox{$n^2$-tuple} obtained by sequentially appending its $n$ columns, and the transpose operation is now understood as a particular involutive $(n^2\times n^2)$-matrix, which we will denote as $\mathcal{T}$.  In this vector format, after dividing both sides by $\text{det}(s)$, \mbox{(\ref{11/5/21.1}) becomes} \begin{equation}\label{6/11/25.2}\mathcal{T}s^{-1} = \frac{\nabla\text{det}(s)}{\text{det}(s)} = \nabla\left(\log \text{det}(s)\right).\end{equation}  Thus,  $\mathcal{T}$ is an uncurling metric for the real matrix algebra $M_n(\mathbb{R})$, as are all scalar multiples \mbox{of $\mathcal{T}$}. 
    

Being an uncurling metric means that integration of $\mathcal{T}s^{-1}$ is path-independent in some open ball centered at the multiplicative identity $\mathbf{1}_{M_n({\mathbb{R})}}$ of $M_n(\mathbb{R})$.  
 Hence, \mbox{$\int_{\mathbf{1}_{M_n({\mathbb{R})}}}^s (\mathcal{\mathcal{T}}t^{-1})\cdot \mathbf{d}t$} is path-independent in the latter neighborhood.  Now consider an $n$-dimensional algebra $\bar{A}$ whose left regular representation $A$ is thus a subalgebra of $M_n(\mathbb{R})$.  Then the above integral is path-independent for paths in the above neighborhood that are additionally restricted to be in (the vector space of elements of) $A$.  Suppose $\mathcal{T}$ maps $A$ to $A$, i.e., $A$ is self-adjoint as a subalgebra of $M_n(\mathbb{R})$. 
 Then the integral is not only path-independent when the domain of integration is restricted to $A$, but on such a path the integrand also remains in $A$.  This means that 
  $d((\mathcal{T}s^{-1})\cdot \mathbf{d}s) = 0$ in some neighborhood of the multiplicative identity $\mathbf{1}_A$ of $A$.  That is,  $\mathcal{T}$ restricted to the elements of $A$ is an uncurling metric for $A$.  Hence, by Theorem \ref{6/6/23.1}, $\mathcal{I}^T\mathcal{T}\mathcal{I}$ is an uncurling metric for $\bar{A}$ (and so are all scalar multiples), where $\mathcal{I}:\bar{A}\rightarrow A$ is the left regular representation isomorphism.  
  
 By famous theorems of Frobenius and Wedderburn, a simple real finite-dimensional unital associative algebra is isomorphic to one of the real algebras $M_q(\mathbb{R})$, $M_q(\mathbb{C})$, or $M_q(\mathbb{H})$, for some positive integer $q$ (where $\mathbb{H}$ is the quaternions).  The left regular representation of each of these algebras is self-adjoint as a subalgebra of $M_k(\mathbb{R})$ for some integer $k$.  Consequently, the prior paragraph indicates that if $\bar{A}$ is simple, it has a space of uncurling metrics that is at least one-dimensional.  
 
 Now suppose the left regular representation of $\bar{A}$ is not self-adjoint as a subalgebra of $M_n(\mathbb{R})$. Then there exists an epimorphism $\Omega:\bar{A}\rightarrow \underline{A}$ such that $ \underline{A}$ {\it is} a simple algebra, and therefore $\underline{A}$ has an uncurling metric $\underline{L}$.  A simple modification of the first three sentences of the proof of Theorem \ref{6/6/23.1} (where we exchange isomorphism $K$ for epimorphism $\Omega$), indicates that $\Omega^T \underline{L}\Omega$ satisfies (\ref{5/22/25.3}) and is thus an uncurling metric for $\bar{A}$ - along with all its scalar multiples.   That is, with respect to the elements $\underline{s}\in \underline{A}$ and the elements $\bar{s}\in\bar{A}$, we consider the linear function $\underline{s}= \Omega \bar{s}$.  Then, $0 = d\left((\underline{s}^{-1})^T \underline{L} \mathbf{d}\underline{s}\right) =  d\left(((\Omega \bar{s})^{-1})^T \underline{L} \mathbf{d}(\Omega \bar{s})\right) = d\left((\bar{s}^{-1})^T \big[\Omega^T \underline{L} \Omega \big]\mathbf{d}\bar{s}\right)$, where we have used Definition \ref{7/21/23.1}, the commutativity of a differential with a linear transformation, and $(\Omega \bar{s})^{-1} = \Omega (\bar{s}^{-1})$ (since $\Omega$ is an epimorphism of rings).   
  \end{proof}
 
 
Before further exploiting Theorem \ref{6/6/23.1}, we  observe that computation of $\mathfrak{u}$ is straightforward in principle.  The algorithm is contained in the proof of the following.  Since it will cause no confusion, in the sequel we simply write an algebra element in the standard basis as, e.g., $s=(x_1,\dots,x_n)$, where the latter is understood to be a column vector.

\begin{proposition}\label{9/2/24.1} The members of $\mathfrak{u}$ are the realizations of a symmetric $(n\times n)$-matrix $M_{\mathfrak{u}}$ whose entries are linear combinations of the components of an $m$-tuple of real parameters, where $m$ is the dimension of $\mathfrak{u}$ as a vector space, and where a realization of $M_{\mathfrak{u}}$ as a real matrix is understood to result from replacing its $m$ parameters with the members of any set of $m$ members of $\mathbb{R}$.   
\end{proposition}

\begin{proof} For an $n$-dimensional algebra with elements $s = (x_1,x_2,\dots,x_n)$ with respect to the standard basis, one takes an initial candidate for $M_\mathfrak{u}$ to be the symmetric $(n\times n)$-matrix $M$ whose upper triangular entries are each a different member of a set of $\frac{n(n+1)}{2}$ real parameters (in effect, the initial candidate is taken to be all of $\text{Sym}_n$).  Computation of $s^{-1}$ produces a vector whose components are each a rational function of the $x_1,x_2,\dots,x_n$.  So far, $Ms^{-1}$ can be understood as a parametrized \mbox{$n$-dimensional} vector field whose components (which can easily be computed) are dependent on the initial set of parameters in $M$ and the $x_1,x_2,\dots,x_n$.  Now, there are $\frac{n(n-1)}{2}$ constraint equations that must be satisfied in order that (\ref{5/22/25.3}) hold for a candidate $L=M$.  Each of these equations has the form, \begin{equation}\label{9/2/24.2} \frac{\partial}{\partial x_i}([Ms^{-1}]_j) - \frac{\partial}{\partial x_j}([Ms^{-1}]_i) = 0,\end{equation} where $[Ms^{-1}]_k$ indicates the $k$-th component of $Ms^{-1}$.  It is easily recognized that in general $M$ will need to be modified in order to satisfy these constraint equations, and this is accomplished as follows.  

As noted, we can write $s^{-1} = (\hat{x}_1,\dots,\hat{x}_n)$, where each $\hat{x}_k$ is a rational function of the components of $s$.  Consequently, $[Ms^{-1}]_k$ is a linear combination of rational functions of the components $\hat{x}_1,\dots,\hat{x}_n$, where the coefficients defining this linear combination are linear combinations of various of the above parameters.  It follows that $\frac{\partial}{\partial x_p}([Ms^{-1}]_q)$ is again a sum rational functions of the components $x_1,\dots,x_n$, where this expression is again linear in the parameters.  Thus, multiplying both sides of (\ref{9/2/24.2}) by a polynomial determined by the polynomial denominators of the constituent rational functions, the left-hand-side of (\ref{9/2/24.2}) can be converted into a polynomial in the components of $s$, while remaining an expression that is linear in the parameters - and the right-hand-side of (\ref{9/2/24.2}) remains zero under this procedure.
In this way, satisfaction of each of the $\frac{n(n-1)}{2}$ constraint equations (\ref{9/2/24.2}) ultimately reduces to the requirement that a particular multivariate polynomial be {\it identically} zero.  After simplification of the polynomial to be a linear combination of independent monomials, each monomial term is then comprised of a coefficient consisting of a linear combination of some of the above $\frac{n(n+1)}{2}$ parameters, multiplying an expression consisting of a product of some of the components of $s$, each raised to some power.  Therefore, each of the above coefficients of these (independent) monomials are required to be zero.  Accordingly, one treats each of the $\frac{n(n-1)}{2}$ constraint equations in the above fashion one after the other in some sequence of steps.  After each individual step, the number of surviving independent parameters (i.e., parameters not required to be equated to some linear combination of other parameters, or set to 0) is in general reduced as the $\frac{n(n-1)}{2}$ constraint equations are successively treated.  $M$ is thereby sequentially modified as necessary as each successive constraint equation is treated.  After treatment of all the constraint equations by executing the above steps in some sequence, the final number $m$ of surviving independent parameters is the dimension of $\mathfrak{u}$, and the final entries of the sequentially modified $M$ evidently yield the requisite $M_\mathfrak{u}$.
\end{proof}

\enlargethispage{2cm}
Consequently, $\mathfrak{u}$ can be represented as a symmetric $(n\times n)$-matrix $M_{\mathfrak{u}}$ whose entries are particular linear combinations of $m$ independent parameters (where $m$ is the dimension of $\mathfrak{u}$), because each uncurling metric results from a single set of selected real values for those parameters.  The determinant of $M_{\mathfrak{u}}$, written as ${\rm det}(M_{\mathfrak{u}})$, is a polynomial in some or all of the components of an ordered $m$-tuple of parameters, $(\sigma_1,\dots,\sigma_m)$.  If $m$ is small  we will often denote the parameters using different Greek letters like $\alpha$, $\beta $, $\gamma$, etc. In the sequel, $\mathfrak{u}$ and $M_{\mathfrak{u}}$ will be used interchangeably depending on the context.


   With the above in mind, it is seen that (\ref{5/23/23.1}) allows us to immediately identify various algebra isomorphism invariants resulting from the anti-rotor. For example,
      
       
  \begin{corollary}\label{8/18/23.2} The following are algebra isomorphism invariants. \begin{itemize} \item The dimension $m$ of $\mathfrak{u}$.   \item The largest rank  of a member of $\mathfrak{u}$.  \item The smallest rank of a nonzero member of $\mathfrak{u}$. \item The number of parameters on which ${\rm det}(M_{\mathfrak{u}})$ depends. \item Invariants of the real algebraic set, \[V_\mathbb{R}({\rm det}(M_{\mathfrak{u}})) = \{(\sigma_1,\dots,\sigma_m)\in\mathbb{R}^m: {\rm det}(M_{\mathfrak{u}}) = 0 \}.\] \end{itemize} \end{corollary}
   
  \section{The dual of the anti-rotor and resulting geometries} 
  \subsection{The transform of an uncurling metric}
  

 Given an uncurling metric $L$, we have that $Ls^{-1}$ is a gradient vector field involving $s^{-1}$.  Thus, we can think of it as deriving from a logarithm.  Consequently, it makes sense to introduce a function $\ell(s)$ satisfying $\ell(\mathbf{1})=1$ and,
  \begin{equation}\label{11/7/20.1} Ls^{-1}  = \nabla\left[\|{\bf 1}\|^2\log  \ell(s)\right] =   \left(\frac{\|{\bf 1}\|^2}{ \ell(s)}\right)\nabla \ell(s),\end{equation} where $\mathbf{1}$ is the algebra's multiplicative identity, and $\|{\bf 1}\|^2 \equiv {\bf 1}\cdot {\bf 1}$. 

   \begin{definition}\label{9/21/23.2} Let $L$ be an uncurling metric for some algebra $A$.  For $s$ in an open ball $D_A$ centered at the multiplicative identity $\mathbf{1}$ and devoid of non-units, with $\|{\bf 1}\|^2 \equiv \mathbf{1}\cdot\mathbf{1}$, the associated {\it unital norm} $ \ell:D_A\rightarrow \mathbb{R}$ is the function determined by, 
  \begin{equation}\label{11/7/20.2} \ell(s) \equiv \exp\left(\frac{1}{\|{\bf 1}\|^2}\int_{\bf 1}^s \left(Lt^{-1}\right)\cdot \mathbf{d}t\right).\end{equation} 
   The set of unital norms associated with an algebra, which is an abelian group $\mathfrak{u}^d_A$ with respect to the usual point-wise operation of multiplication of real functions, is called the {\it unital norm group}, also known as the {\it dual of the anti-rotor}. $\mathfrak{u}^d_A$ may also be denoted as $\mathfrak{u}^d$ if no confusion arises.  
   \end{definition} 
   
The integral in (\ref{11/7/20.2}) is path-independent due to (\ref{5/22/25.3}).  

The duality referred to in Definition \ref{9/21/23.2} is easily appreciated in its relation to integral transforms.  Equation (\ref{11/7/20.2}) can be rewritten as  \mbox{$\log\ell(s) = \frac{1}{\|\mathbf{1}\|^2}\int_{{\bf 1}}^{s} \left[L t^{-1}\right]\cdot \mathbf{d}t$.}   We can view this as a (forward) integral transform, mapping $Ls^{-1}$ to $\log \ell(s)$.  But, based on the assumed standard basis, an expression for $s^{-1}$ is known for any given algebra (i.e., matrix $L$ can be extracted from $Ls^{-1}$, and $Ls^{-1}$ can be calculated from $L$), so this effectively results in a ``forward transform" of $L$ to $\ell(s)$ (of course, this transform is restricted to the space of uncurling metrics).    Equation (\ref{11/7/20.2}) also implies (\ref{11/7/20.1}), the latter being an (inverse) integral transform from $\log \ell(s)$ to $Ls^{-1}$ in the sense of distributions.  Thus we have an ``inverse transform" of $\ell(s)$ to $L$ (restricted to the space of unital norms).
This duality is specifically defined with respect to the multiplicative inversion operation $s\mapsto s^{-1}$.
  
 
 Based on the mapping from $\mathfrak{u}$ to $\mathfrak{u}^d$ defined by (\ref{11/7/20.2}),  we can write \begin{equation}  \mathfrak{u}^d = \exp\left(\int_{\bf 1}^s [\mathfrak{u}t^{-1}]\cdot \mathbf{d}t\right). \label{6/7/23.1}\end{equation}  
 In the above, the notation $\mathfrak{u}t^{-1}$ means $\{Lt^{-1}: L\in \mathfrak{u}\}$.  Accordingly, the right-hand-side of (\ref{6/7/23.1}) is understood to be a collection of exponentiated indefinite integrals.   Note that in (\ref{6/7/23.1}) we do not need the factor $\frac{1}{\|{\bf 1}\|^2}$, since $\frac{1}{\|{\bf 1}\|^2}\mathfrak{u}=\mathfrak{u}$.  The entities  $\mathfrak{u}$ and $\mathfrak{u}^d$ are considered dual based on the invertible transform relation (\ref{11/7/20.2}).
  

\subsection{Special unital norms}
Of particular interest are those uncurling metrics such that elements that are multiplicative inverses of each other also behave inversely with respect to that metric.   
 
 \enlargethispage{.5cm}
    \begin{definition}\label{9/20/23.1} For a given algebra, a {\it normalized uncurling metric} $L_{\shortparallel}$ is an uncurling metric that satisfies, \begin{equation} \label{7/7/21.1} s^TL_{\shortparallel} s^{-1} =\|{\bf 1}\|^2,\end{equation} for $s,s^{-1}$ in an open ball centered at $\mathbf{1}$ devoid of non-units.   A {\it special unital norm} $\ell_{\rm sp}(s)$ is a unital norm resulting from (\ref{11/7/20.2}) when $L$ is a normalized \mbox{uncurling metric $L_{\shortparallel}$}. 
\end{definition}

The special unital norms arising from normalized uncurling metrics have the following useful properties.
  
  \begin{proposition}\label{9/24/23.1} $ \ell_{\rm sp}(s)$ is a degree-1 positive homogeneous function, and
\begin{equation}\label{5/23/23.3}  \ell_{\rm sp}(s^{-1}) = \big(\ell_{\rm sp}(s)\big)^{-1} = \frac{1}{\ell_{\rm sp}(s)}, \end{equation} when $s,s^{-1}$ are in an open ball centered at $\mathbf{1}$ devoid of non-units. \end{proposition}

\begin{proof}  The requisite open ball in the theorem statement always exists, as was already observed immediately following the statement of Definition \ref{7/21/23.1}.  Let domain $D$ be such a neighborhood of $\mathbf{1}$.

Given $\alpha>0$ such that $\frac{1}{\alpha}s$ and $\alpha s$ are both in $D$, (\ref{11/7/20.2}) implies,
\begin{eqnarray}\nonumber  \log\ell_{\rm sp}(\alpha s) &=&  \frac{1}{\|{\bf 1}\|^2}\int_{\bf 1}^{\alpha s} \left(L_{\shortparallel}t^{-1}\right)\cdot \mathbf{d}t = \frac{1}{\|{\bf 1}\|^2}\int_{\frac{\mathbf{1}}{\alpha}\mathbf{1}}^{s} \left(L_{\shortparallel}(\alpha w)^{-1}\right)\cdot \mathbf{d}(\alpha w)  \\ \nonumber &=& \frac{1}{\|{\bf 1}\|^2}\int_{\frac{\mathbf{1}}{\alpha} \mathbf{1}}^{\mathbf{1}} \left(L_{\shortparallel}w^{-1}\right)\cdot \mathbf{d}w  +  \frac{1}{\|{\bf 1}\|^2}\int_{\mathbf{1}}^{s} \left(L_{\shortparallel}w^{-1}\right)\cdot \mathbf{d}w  \\ &=&   \log\alpha + \log \ell_{\rm sp}(s), \label{4/2/23.1}
\end{eqnarray} where the integral from $\frac{\mathbf{1}}{\alpha} \mathbf{1}$ to $\mathbf{1}$ can be easily evaluated along the line segment with those endpoints, and we have used (\ref{7/7/21.1}) in evaluation of that integral to  obtain the term $\log \alpha$ on the right-hand-side of the final equality.  Thus, we have $\ell_{\rm sp}(\alpha s) = \alpha\ell_{\rm sp}(s)$, degree-1 positive homogeneity.


Define $D^{-1}$ to be the set consisting of the multiplicative inverses of all members of $D$.  It is easily shown that $D^{-1}$ is an open neighborhood of $\mathbf{1}$, and the domains \mbox{$D$, $D^{-1}$}, are diffeomorphic via the multiplicative inversion operation.  
We now consider a smooth path $\mathcal{P}$ from $\mathbf{1}$ to $s$ which is inside $D$.  Let $\mathcal{P}^{-1}$ be the point set consisting of the multiplicative inverses of the members of $\mathcal{P}$.  Clearly, $\mathcal{P}^{-1}$ is a smooth path contained in $D^{-1}$.  
 We then have, 
\begin{eqnarray}  \log \ell_{\rm sp}(s^{-1})   &=&   \frac{1}{\|{\bf 1}\|^2} \int_{\bf 1}^{s^{-1}} \left(L_{\shortparallel}t^{-1}\right) \cdot \mathbf{d}t \nonumber \\ &=&  \frac{1}{\|{\bf 1}\|^2} \int_{\bf 1}^s \left(L_{\shortparallel}y\right)\cdot \mathbf{d}\left(y^{-1}\right) \nonumber \\ &=&  \frac{1}{\|{\bf 1}\|^2}\left(\left(y^{-1}\cdot \left(L_{\shortparallel}y\right)\right)\bigg|_{\bf 1}^s - \int_{\bf 1}^s y^{-1}\cdot \mathbf{d}\left(L_{\shortparallel}y\right)\right)   \nonumber \\ &=& -\frac{1}{\|{\bf 1}\|^2} \int_{\bf 1}^s y^{-1}\cdot (L_{\shortparallel}\mathbf{d}y) = -\frac{1}{\|{\bf 1}\|^2} \int_{\bf 1}^s \left(L_{\shortparallel}y^{-1}\right) \cdot \mathbf{d}y  \nonumber \\ &=& -\log \ell_{\rm sp}(s),\label{4/2/23.5} \end{eqnarray} where we have used the change of variable $y=t^{-1}$, commutativity of a linear transformation and a differential, the property that $L_{\shortparallel}$ is a real symmetric matrix (i.e., self-adjoint), and  (\ref{7/7/21.1}).  Equation (\ref{5/23/23.3}) then follows.
\end{proof}

Using an argument similar to that employed in the second paragraph of the proof of the above proposition, it can also be shown more generally that  \[  \ell(\alpha s) = \alpha^\gamma \ell(s),\] with $\gamma = \frac{\mathbf{1}^T L \mathbf{1}}{\|\mathbf{1}\|^2}$, so long as there is an open ball centered at $\mathbf{1}$ devoid of non-units but containing the line segment from $\frac{1}{\alpha}\mathbf{1}$  to $\alpha\mathbf{1}$.  Thus, even a ``non-special" $ \ell(s)$ is a positive homogeneous function, albeit of degree $\gamma$.  


As shown in Section 4.4, the argument in the proof of Proposition \ref{6/11/25.1} can be augmented to show that there is always at least one solution to (\ref{5/22/25.3}), (\ref{7/7/21.1}), i.e.,  a normalized uncurling metric always exists.  This derives from the fact that $\mathcal{T}$ is actually a normalized uncurling metric for $M_n(\mathbb{R})$.  The latter is seen by taking the dot product with $s$ on both sides of the first equality in (\ref{6/11/25.2}) to obtain
\begin{equation}\label{9/15/25.1}s\cdot\mathcal{T}s^{-1} = \frac{s\cdot\nabla\text{det}(s)}{\text{det}(s)} = n = \|\mathbf{1}_{M_n(\mathbb{R})}\|^2,\end{equation} 
where we have invoked the Euler Homogeneous Function Theorem as applied to $\text{det}(s)$.    Thus, (\ref{7/7/21.1}) is satisfied.   

As also formalized in Section 4.4, this program additionally generalizes the ``usual" algebra norm (the determinant of an element's left regular representation, e.g., \cite{bourbaki:1989}), since the latter is closely related to one of the $\ell_{\rm sp}(s)$. Keeping in mind details of the proof of Proposition \ref{6/11/25.1},  this can begin to be appreciated by observing that if the algebra in question is $M_n(\mathbb{R})$ (for which $\|\mathbf{1}_{M_n(\mathbb{R})}\|^2=n$), and $L$ is taken to be the matrix transpose operation $\mathcal{T}$ (which is a normalized uncurling metric as above), then (\ref{6/11/25.2}) indicates that the corresponding solution to (\ref{11/7/20.1}) is given by $\ell_{\rm sp}(s) = (\text{det}(s))^\frac{1}{n}$. 

\subsection{The geometries associated with an algebra}

As discussed in \cite{fgreensite:2025b}, George Birkhoff showed that his four-postulate formulation of Euclidean geometry implies that the plane is an affine space whose vector space component is the Cartesian plane \cite{birkhoff:1932}.  Despite the fact that his formulation is in principle analytic (since it intrinsically incorporates $\mathbb{R}$), as preparation he provided what is essentially a synthetic proof of the Pythagorean Theorem.  Ultimately, according to his program, the theorem is equivalent to the statement that $\mathbb{R}\oplus\mathbb{R}$ is endowed with a length function satisfying $\ell(x,y)=\sqrt{x^2+y^2}$.  As observed in \cite{fgreensite:2025b}, for $s=(x,y)\ne (0,0)$, it is a simple matter to demonstrate that the equation in the prior sentence is itself equivalent to, \begin{eqnarray*} s = \ell(s) \nabla \ell(s),\\  \ell(\nabla \ell(s)) = 1. \end{eqnarray*}  So we could instead take the above characterization of the non-origin points of $\mathbb{R}\oplus\mathbb{R}$ to be the Pythagorean Theorem.  If one considers the {\it format} of this characterization to be the important feature, we might deem the following (and the suggested appellation) as faithful to its essence.

\begin{corollary}[The Pythagorean Theorem for algebras]\label{10/7/25.1} For a given algebra, suppose $L_{\shortparallel}$ is a nonsingular normalized uncurling metric and $\ell_{\rm sp}(s)$ is its dual.  Then for elements $s$ in a unit ball around $\mathbf{1}$ devoid of non-units,
   \begin{eqnarray}\label{5/22/24.5}  s = \ell_{\rm sp}(s)\Big(\|\mathbf{1}\|^2 L_{\shortparallel}^{-1}[\nabla\ell_{\rm sp}(s^{-1})]\Big),\\
\label{5/22/24.6}\ell_{\rm sp}\Big(\|\mathbf{1}\|^2 L_{\shortparallel}^{-1}[\nabla\ell_{\rm sp}(s^{-1})]\Big) = 1,\,\end{eqnarray}
where $\nabla\ell_{\rm sp}(s^{-1})$ means that $\nabla\ell_{\rm sp}$ is evaluated at $s^{-1}$.
\end{corollary} 
\begin{proof}
According to Proposition \ref{9/24/23.1}, $\frac{1}{\ell_{\rm sp}(s)}  = \ell_{\rm sp}(s^{-1})$.  Applying this to the right-hand-side of (\ref{11/7/20.1}), next applying $L_{\shortparallel}^{-1}$ to both sides of the resulting equation, and then substituting $s^{-1}$ for $s$, we obtain (\ref{5/22/24.5}).  Now applying $\ell_{\rm sp}$ to both sides of the latter equation, and invoking the degree-1 positive homogeneity of $\ell_{\rm sp}$ that is another consequence of Proposition \ref{9/24/23.1}, we obtain (\ref{5/22/24.6}).
\end{proof}
  

The equation set in Corollary \ref{10/7/25.1} has an even more obvious similarity to an equation set that holds not just for Euclidean space, but for quadratic spaces in general. 
Thus, if $L$ is the scalar product matrix associated with a quadratic space where the associated quadratic form is positive-definite or indefinite but nondegenerate, it follows that
  \begin{eqnarray}\label{5/22/25.1} s = \ell_{\rm qs}(s) L^{-1}[\nabla \ell_{\rm qs}(s)],\\ \ell\left(L^{-1}[\nabla \ell_{\rm qs}(s)]\right) = 1, \label{5/22/25.11}\end{eqnarray} where the above equations pertain to a domain in $\mathbb{R}^n$ on which the quadratic form is strictly positive and written as $\ell_{\rm qs}^2(s)$, and $\ell_{\rm qs}(s)$ is the positive square-root.  To see this, begin with the quadratic form defined by $ \ell_{\rm qs}^2(s) \equiv s^TLs$.  Take the gradient of both sides of the latter equation and then apply $L^{-1}$ to both sides of the resulting equation to obtain (\ref{5/22/25.1}).  
  Since a quadratic form is a degree-2 positive homogeneous function, it follows that $\ell_{\rm qs}(s)$ is degree-1 positive homogeneous.  Now apply $\ell_{\rm qs}$ to both sides of  (\ref{5/22/25.1}) and invoke the homogeneity of $\ell_{\rm qs}$ to obtain (\ref{5/22/25.11}).      
Also, (\ref{5/22/25.1}) implies \begin{equation}\label{5/26/25.1} Ls = \nabla\left(\frac{1}{2}\ell_{\rm qs}^2(s)\right).\end{equation}  The exponent on the right-hand-side indicates the reason for the name ``quadratic space".  Comparison of (\ref{5/26/25.1}) with the first equality in (\ref{11/7/20.1}) justifies referring to the algebra-associated space described by (\ref{5/22/24.5}), (\ref{5/22/24.6}) as a ``logarithmic space".

Now, from (\ref{5/22/25.1}), {\it in a quadratic space the square of the norm} $\ell_{\rm qs}(s)$ is obtained from path-independent integration of $t$ with respect to the metric (symmetric tensor) $L_{\rm sym}$ according to, \begin{equation}\label{8/19/25.1} \ell_{\rm qs}^2(s) = 2\int_0^s (L_{\rm sym}t)\cdot \mathbf{d}t =  2\int_0^s t^T L_{\rm sym} \mathbf{d}t.\end{equation}
From (\ref{5/22/25.1}), (\ref{5/22/25.11}), we are interested in the associated unit sphere as the locus of points satisfying $\ell_{\rm qs}(s)=1$ since (as typified with Eucldean space or Minkowski space) the angle associated with two vectors $s_1$, $s_2$ is defined as the arclength of the geodesic path on the unit sphere connecting $\frac{s_1}{\ell_{\rm qs}(s_1)}$, $\frac{s_2}{\ell_{\rm qs}(s_2)}$ where, as with the integration on the right-hand-side of (\ref{8/19/25.1}) that is taken with respect to the metric $L_{\rm sym}$\,, arclength and geodesic are also defined with respect to the metric $L_{\rm sym}$\,. 

Analogously, from (\ref{11/7/20.1}), {\it in a logarithmic space the logarithm of the norm} $\ell_{\rm sp}(s)$ is obtained from path-independent integration of $t^{-1}$ with respect to the uncurling metric $L_{\shortparallel}$ \mbox{according to},
\begin{equation}\label{8/19/25.2} \log\ell_{\rm sp}(s) =  \frac{1}{\|{\bf 1}\|^2} \int_{\bf 1}^s \left(L_{\shortparallel}t^{-1}\right) \cdot \mathbf{d}t =
\frac{1}{\|{\bf 1}\|^2} \int_{\bf 1}^s \left( t^{-1}\right)^T L_{\shortparallel} \mathbf{d}t .\end{equation} 
 From
(\ref{5/22/24.5}), (\ref{5/22/24.6}), the unit sphere is now the locus of points satisfying $\ell_{\rm sp}(s) = 1$.    The angle between two vectors $s_1$, $s_2$ is defined in the same manner as accomplished in a quadratic space, i.e., as the arclength of the geodesic path on the unit sphere connecting $\frac{s_1}{\ell_{\rm sp}(s_1)}$, $\frac{s_2}{\ell_{\rm sp}(s_2)}$ where, as with the integration on the right-hand-side of (\ref{8/19/25.2}) that is taken with respect to the uncurling metric $L_\shortparallel$\,, arclength and geodesic are also defined with respect to the uncurling \mbox{metric $L_{\shortparallel}$}\,.

Thus, as with the distinct geometries arising from different signatures of scalar products defining quadratic spaces, distinct geometries associated with an algebra (in the sense of assigned ``norms" and ``angles") will result from different classes of normalized uncurling metrics - defining the logarithmic spaces.

However, any uncurling metric $L$ results in ``norms" and ``angles" in the same way as the latter are obtained with $L_\shortparallel$.
The following shows that the implied collection of geometries associated with an algebra $A$ (i.e., resulting from the members of $\mathfrak{u}_A^d$ as above) persists under an algebra isomorphism.  

  \begin{theorem}\label{6/6/23.2} 
If $A_1$ and $A_2$ are isomorphic algebras, then the respective collections of geometries implied by $\mathfrak{u}_{A_1}^d$ and  $\mathfrak{u}_{A_2}^d$ are the same.
 \end{theorem}

\begin{proof} Let $K:A_1\rightarrow A_2$ be the algebra isomorphism, where for $s_1\in A_1$ we define $s_2\equiv Ks_1$. 
Equation (\ref{5/23/23.1}) implies,
\begin{eqnarray}\nonumber \int_{{\bf 1}_{A_1}}^{s_1} \left(\mathfrak{u}_{A_1}t_1^{-1}\right)\cdot \mathbf{d}t_1 &=& \int_{{\bf 1}_{A_1}}^{s_1} \left( K^T\mathfrak{u}_{A_2}Kt_1^{-1}\right)\cdot \mathbf{d}t_1 \\ \nonumber &=& 
\int_{{\bf 1}_{A_2}}^{s_2} \left( K^T\mathfrak{u}_{A_2}K (K^{-1}t_2)^{-1}\right)\cdot \mathbf{d}\left(K^{-1}t_2\right) \\ &=& \int_{{\bf 1}_{A_2}}^{s_2} \left(K^T\mathfrak{u}_{A_2}t_2^{-1}\right)\cdot \left(K^{-1}\mathbf{d}t_2\right) = \int_{{\bf 1}_{A_2}}^{s_2} \left(\mathfrak{u}_{A_2}t_2^{-1}\right)\cdot \mathbf{d}t_2, \label{9/6/24.2}
\end{eqnarray} 
where we have used the isomorphism assumption (which implies $(K^{-1}t_2)^{-1} = K^{-1}t_2^{-1}$), the commutativity of a linear transformation and a differential, and exploitation of an adjoint in the context of an inner product.  Recalling (\ref{6/7/23.1}), the theorem follows.
\end{proof}

Obviously, the theorem does require the isomorphism assumption.  An arbitrary basis change applied to $A_1$ would not in general result in the dual of the anti-rotor of some other algebra $A_3$ such that $A_3$ is isomorphic to $A_1$.

\subsection{Generalization of the usual algebra norm of an element}
As already noted, the usual algebra norm is given by the determinant of an element's left regular representation.  The dual of the anti-rotor (as the group of unital norms) subsumes this definition in the sense of Theorem \ref{11/17/22.1}, as derived in this subsection.  
  
 \begin{proposition} \label{4/8/22.1} $\mathfrak{u}_{M_n(\mathbb{R})}$ is one-dimensional and there is only a single normalized uncurling metric, this being the  transpose operation on $(n\times n)$-matrices expressed as an $(n^2\times n^2)$-matrix. The associated special unital norm is $({\rm det}(s))^{\frac{1}{n}}$, i.e., the usual norm raised to the power given by the reciprocal of the dimension of $M_n(\mathbb{R})$. \end{proposition}
 \begin{proof} The vector space of members of $ M_n(\mathbb{R})$ consists of  vectors with $n^2$ components.   Using the sequentially appended matrix columns convention as introduced in the proof of Proposition \ref{6/11/25.1}, we can equivalently represent a member $s\in M_n(\mathbb{R})$ in either matrix or column vector format.   An uncurling metric for this algebra will be a symmetric matrix $L\in M_{n^2}(\mathbb{R})$ that must satisfy (\ref{5/22/25.3}), i.e., $d\left([Ls^{-1}]\cdot \mathbf{d}s\right)=0$. This constitutes a condition on every pair of distinct components of $Ls^{-1}$ (the latter again having $n^2$ components).  Thus,  the partial derivative of the  $(i,j)$-entry of matrix $Ls^{-1}$ (say), taken with respect to the $(k,m)$-entry of matrix $s$ (say) (with $(k,m)\ne (i,j)$),  minus  the partial derivative of the $(k,m)$-entry of $Ls^{-1}$  taken with respect to the $(i,j)$-entry of matrix $s$, must be zero.
 
 To examine this constraint, let $s_{pq}$ be the $(p,q)$-entry of matrix $s$.   Differentiating both sides of the expression
$ ss^{-1} = {\bf 1}$ with respect to $s_{pq}$ (as in Matrix Calculus), we obtain $\frac{\partial (s^{-1}s)}{\partial s_{pq}} = \frac{\partial (s^{-1})}{\partial s_{pq}}s + s^{-1}\frac{\partial s}{\partial s_{pq}}=0$ (where $0$ is the matrix of all zero entries).  Thus, we have the well-known result 
 \[ \frac{\partial (s^{-1})}{\partial s_{pq}} = -s^{-1}\frac{\partial s}{\partial s_{pq}}s^{-1} = -s^{-1}e_{pq}s^{-1},\]   where $e_{pq}$ is the matrix whose $(p,q)$ entry is $1$, but with all other entries zero (and $\{ e_{pq}\}$ is the canonical basis for $ M_n(\mathbb{R})$).  
 Hence, the partial derivative of $s^{-1}$ with respect to the  $(k,m)$-entry of matrix $s$, which is itself a matrix, (``the derivative of $s^{-1}$ in the $e_{km}$-direction") is, \begin{equation}\label{8/18/22.5} -s^{-1} e_{km}s^{-1} = -(s^{-1})_{:k} (s^{-1})_{m:},\end{equation} where $(s^{-1})_{:k}$ indicates the $k$-th column of $s^{-1}$, and $(s^{-1})_{m:}$ indicates the $m$-th row of $s^{-1}$.    Similarly, the derivative of $s^{-1}$ with respect to the $(i,j)$-entry of $s$ (``the derivative of $s^{-1}$ in the $e_{ij}$-direction")  is, \begin{equation}\label{8/18/22.6} -s^{-1} e_{ij}s^{-1} = -(s^{-1})_{:i} (s^{-1})_{j:}.\end{equation}  It is thus clear that the sets of entries of the $e_{km}$-direction derivative matrix and the $e_{ij}$-direction derivative  matrix share the member  $-(s^{-1})_{mi} (s^{-1})_{jk}$, and they can share no other member because $\{(s^{-1})_{pq}\} = \{\frac{C_{qp}}{\det(s)}\}$, where the right-hand-side of this equation is clearly a set of linearly independent functions since $\{C_{qp}\}$ is the matrix of cofactors of $s$.  The latter observation also indicates that the union of the entries of the $e_{km}$-direction derivative matrix with the entries of the $e_{ij}$-direction derivative matrix is a set of linearly independent functions.

Therefore, since a derivative and a linear transformation commute, the partial derivative in the $e_{km}$-direction of the $e_{ij}$-component of $Ls^{-1}$ is a linear combination of the $n^2$ entries of the matrix on the right-hand-side of (\ref{8/18/22.5}), and  the partial derivative in the $e_{ij}$-direction of the $e_{km}$-component of $Ls^{-1}$ is a linear combination of the $n^2$ entries on the right-hand-side of (\ref{8/18/22.6}).
So satisfaction of \mbox{$d\left([Ls^{-1}]\cdot \mathbf{d}s\right) = 0$,} based on the observation made at the end of the prior paragraph, requires that all entries in any two rows of $L$ must be zero, except for the one entry in each row corresponding to the shared member of the two relevant partial derivative matrices (the right-hand-sides of (\ref{8/18/22.5}) and (\ref{8/18/22.6})) noted in the prior paragraph - and these two corresponding nonzero entries of the two cited rows of $L$ must have the same value.  This prescription must be true for any two rows of $L$, so all nonzero entries of $L$ must have the same value, and there can be at most one nonzero value in each row of $L$.  Thus, the anti-rotor is at most one-dimensional. If it is one-dimensional, the anti-rotor can be represented as a parameter $\alpha$ ranging over $\mathbb{R}$ multiplying a nonzero constant matrix whose nonzero entries are all 1.  
For the special unital norm, the value of $\alpha$ is determined by the requirement that (\ref{7/7/21.1}) be satisfied. 


Hence, there is at most one possible normalized uncurling metric.  But according to observations made in the penultimate paragraph of Section 4.2, $ M_n(\mathbb{R})$ does indeed have a normalized uncurling metric - this being the $(n^2\times n^2)$-matrix $\mathcal{T}$ corresponding to the matrix transpose operation on $(n\times n)$-matrices.  In the final paragraph of Section 4.2 it was shown that $(\mbox{det}(s))^{\frac{1}{n}}$ is the associated special unital norm.  Since the usual norm of an element $s\in M_n(\mathbb{R})$ is $(\mbox{det}(s))^n$, the theorem follows.
 \end{proof}
 
 Expanding on Proposition \ref{4/8/22.1}, the following lemma and theorem indicate that the unital norms are an elaboration of the usual algebra norm.  Recall that a subalgebra of $M_n(\mathbb{R})$ is defined to be self-adjoint if the transpose operation maps the subalgebra to itself.
 

 
  \begin{lemma}\label{5/13/22.1} If the left regular representation of an $n$-dimensional algebra $\bar{A}$ with elements $\bar{s}$ is self-adjoint as a subalgebra of $M_n(\mathbb{R})$, then there is a special unital norm given by the usual algebra norm raised to the power of $\frac{1}{n}$, i.e., $\ell_{\rm sp}(\bar{s}) = [{\rm det}(\mathcal{I}\, \bar{s})]^\frac{1}{n}$, where $\mathcal{I}$ is the isomorphism from $\bar{A}$ to its left regular representation.
\end{lemma}

\begin{proof} 
Let $\bar{A}$ be an $n$-dimensional algebra and let $A$ be its left regular representation.  Given $\bar{s}\in\bar{A}$ we denote its image under the left regular representation as $s = \mathcal{I}\bar{s}$.
 So, $\bar{A}$ is isomorphic to $A$, and $A$ is a subalgebra of $ M_n(\mathbb{R})$.   $\mathbf{1}_{M_n(\mathbb{R})}$ is the $(n\times n)$-identity matrix, and observe that as an element of $ M_n(\mathbb{R})$, $\mathbf{1}_A$ is likewise the $(n\times n)$-identity matrix.  However, $ M_n(\mathbb{R})$ is $n^2$-dimensional with $\|\mathbf{1}_{M_n(\mathbb{R})}\|^2 = n$, while $A$ is $n$-dimensional with $\|\mathbf{1}_A\|^2 = \|\mathbf{1}_{\bar{A}}\|^2$ (because $\mathbf{1}_{\bar{A}}$ and  $\mathbf{1}_A$ can be represented by the same $n$-tuple in the standard basis).   

As noted in the first paragraph of the proof of Proposition \ref{6/11/25.1}, the transpose operation on $ M_n(\mathbb{R})$ can be expressed as a linear transformation $\mathcal{\mathcal{T}}:\mathbb{R}^{n^2}\rightarrow \mathbb{R}^{n^2}$.  As further noted in the penultimate paragraph of Section 4.2, $\mathcal{T}$ is a normalized uncurling metric associated with algebra  $ M_n(\mathbb{R})$.  Being an uncurling metric means that integration of $\mathcal{\mathcal{T}}s^{-1}$ is path-independent in an open ball devoid of non-units centered at the multiplicative identity of $ M_n(\mathbb{R})$.  
 Hence, $\int_{\mathbf{1}_{M_n({\mathbb{R})}}}^s (\mathcal{\mathcal{T}}t^{-1})\cdot \mathbf{d}t$ is path-independent in the latter neighborhood and, in particular, is path-independent for paths in that neighborhood restricted to be in $A$.  By assumption, $\mathcal{\mathcal{T}}$ maps $A$ to $A$. 
 Thus, $\int_\mathbf{1}^s \mathcal{\mathcal{T}}t^{-1}\cdot \mathbf{d}t$ is path-independent when the domain of integration is restricted to $A$, and on such a path the integrand also remains in $A$.  This means that 
  $d((\mathcal{\mathcal{T}}s^{-1})\cdot \mathbf{d}s) = 0$ in some neighborhood of $\mathbf{1}_A$ in $A$.
  
   Now, we can use $\mathcal{\mathcal{T}}$ to define a scalar product on $A$ (not necessarily positive-definite or nondegenerate) as $\langle s_1,s_2\rangle_A \equiv  \frac{\|\mathbf{1}_{A}\|^2}{n}s_1^T \mathcal{\mathcal{T}} s_2$. 
    Since $A$ and $\bar{A}$ are isomorphic, the above scalar product on $A$ can be employed to produce a scalar product on $\bar{A}$ so that the two resulting quadratic spaces are isomorphic (i.e., with scalar products that are invariant with respect to the application of $\mathcal{I}$).  Specifically, for $\bar{s_1},\bar{s_2}\in\bar{A}$, define $s_1 \equiv \mathcal{I}\bar{s}_1$ and $s_2 \equiv \mathcal{I}\bar{s}_2$.  Then define the induced scalar product on $\bar{A}$ via, \begin{equation}\label{7/15/22.1} \bar{\mathcal{\mathcal{T}}} \equiv  \frac{\|\mathbf{1}_{\bar{A}}\|^2}{n}\mathcal{I}^T\mathcal{\mathcal{T}}\mathcal{I},\end{equation} so that $\langle \bar{s}_1,\bar{s}_2\rangle_{\bar{A}} \equiv \bar{s}_1^T \bar{\mathcal{\mathcal{T}}} \bar{s}_2$.  Applying (\ref{7/15/22.1}) and the previously noted $\|\mathbf{1}_A\|^2 = \|\mathbf{1}_{\bar{A}}\|^2$, we have \begin{equation}\label{8/25/22.1} \langle \bar{s}_1,\bar{s}_2\rangle_{\bar{A}} = \bar{s}_1^T \bar{\mathcal{T}} \bar{s}_2 = \frac{\|\mathbf{1}_{\bar{A}}\|^2}{n}( s_1^T \mathcal{T} s_2) = \langle s_1,s_2\rangle_{A}.\end{equation}  The fraction $\frac{\|\mathbf{1}_{\bar{A}}\|^2}{n}$ is employed in (\ref{7/15/22.1}) in order that $\langle \mathbf{1}_{\bar{A}},\mathbf{1}_{\bar{A}}\rangle_{\bar{A}} = \|\mathbf{1}_{\bar{A}}\|^2$, given that $\mathbf{1}_{M_n(\mathbb{R})}^T \mathcal{T}\,\mathbf{1}_{M_n(\mathbb{R})} = n = \|\mathbf{1}_{M_n(\mathbb{R})}\|^2$. Due to this isomorphism of $A$ and $\bar{A}$ as quadratic spaces (as established in (\ref{7/15/22.1})), the previously demonstrated \mbox{$d((\mathcal{T}s^{-1})\cdot \mathbf{d}s) = 0$} in some neighborhood of $\mathbf{1}_A$ in $A$ implies 
\mbox{$d((\bar{\mathcal{T}}\bar{s}^{-1})\cdot d\bar{s}) = 0$.} Thus, $\bar{\mathcal{T}}$ is an uncurling metric for $\bar{A}$.


In the context of $M_n(\mathbb{R})$, we know from Proposition \ref{4/8/22.1} that $L=\mathcal{T}$ satisfies (\ref{7/7/21.1}), where $\|\mathbf{1}\|^2$ in  (\ref{7/7/21.1}) is $\|\mathbf{1}_{M_n({\mathbb{R})}}\|^2 = n$.  Combining these facts with (\ref{9/15/25.1}) and (\ref{8/25/22.1}), we have \begin{equation}\label{9/27/23.1} \bar{s}^T \bar{\mathcal{T}} \bar{s}^{-1} = \frac{\|\mathbf{1}_{\bar{A}}\|^2}{n}( s^T \mathcal{T} s^{-1}) = \|\mathbf{1}_{\bar{A}}\|^2 .\end{equation} 
Thus, $\mathcal{\bar{T}}$ is a normalized uncurling metric for $\bar{A}$.  Hence, a special unital norm on $\bar{A}$ results from (\ref{11/7/20.2}) (with $\mathbf{1}=\mathbf{1}_{\bar{A}}$)  via the normalized uncurling metric $L_{\shortparallel} = \bar{\mathcal{T}}$.  
  We then obtain, 
\begin{eqnarray}\nonumber \log\ell_{\rm sp}(\bar{s}) = \frac{1}{\|\mathbf{1}_{\bar{A}}\|^2} \int_{\mathbf{1}_{\bar{A}}}^{\bar{s}} (\bar{\mathcal{T}}\bar{s}^{-1})\cdot \mathbf{d}\bar{s}
 &=& \frac{1}{\|\mathbf{1}_{\bar{A}}\|^2} \int_{\mathbf{1}_{A}}^{s} \frac{\|\mathbf{1}_{\bar{A}}\|^2}{n}(\mathcal{T}s^{-1})\cdot \mathbf{d}s \\ &=& \frac{1}{n} \log\text{det}(s),\label{9/27/23.2}\end{eqnarray} where the first equality is from (\ref{11/7/20.2}), the second equality follows from application of (\ref{7/15/22.1}), and the final equation results from Proposition \ref{4/8/22.1}.
The lemma then follows. 
 \end{proof} 
 

So now we consider the general case.
 
  \begin{theorem}\label{11/17/22.1}
  Given any algebra $\bar{A}$, suppose there is an epimorphism $\Omega:\bar{A}\rightarrow \underline{A}$ where $ \underline{A}$ has dimension $m\ge 1$ and such that the left regular representation of $\underline{A}$ is self-adjoint as a subalgebra of $M_m(\mathbb{R})$.  Then there is a special unital norm member of $\mathfrak{u}_{\bar{A}}$ given by $[{\rm det}(\underline{\mathcal{I}}\Omega \bar{s})]^{\frac{1}{m}}$, for $\bar{s}\in\bar{A}$, where $\underline{\mathcal{I}}$ is the isomorphism from $\underline{A}$ to its left regular representation.  Furthermore, at least one such epimorphism $\Omega$ exists. 
   \end{theorem}
 
  \begin{proof}  
  If $\Omega$ exists, then by Lemma \ref{5/13/22.1} there exists a normalized uncurling metric $\underline{L}_{\shortparallel}$ for $\underline{A}$.  For elements $\underline{s}\in \underline{A}$ in a neighborhood of $\mathbf{1}_{\underline{A}}$ devoid of non-units, the associated special unital norm is, \begin{equation}\label{1/9/25.1} \ell_{\rm sp}(\underline{s}) =    \exp\left(\frac{1}{\|\mathbf{1}_{\underline{A}}\|^2}\int_{\mathbf{1}_{\underline{A}}}^{\underline{s}} [\underline{L}_{\shortparallel}\,\underline{t}^{-1}]\cdot \mathbf{d}\underline{t}\right) = [\text{det}(\underline{\mathcal{I}}\, \underline{s})]^\frac{1}{m}.\end{equation}
  We know from the argument in the last paragraph of the proof of Proposition \ref{6/11/25.1} that $\Omega^T \underline{L}_{\shortparallel} \Omega$ is an uncurling metric for $\bar{A}$.  Thus, it is easy to see that we then have a normalized uncurling metric for $\bar{A}$ given by, 
  \begin{equation}\label{1/9/25.2} \bar{L}_{\shortparallel} \equiv \frac{\|\mathbf{1}_{\bar{A}}\|^2}{\|\mathbf{1}_{\underline{A}}\|^2} \Omega^T \underline{L}_{\shortparallel} \Omega, \end{equation} since (\ref{7/7/21.1}) is satisfied because $\Omega\bar{s} =  \underline{s}$ and  $\Omega\bar{s}^{-1} =  \underline{s}^{-1}$ (because $\Omega$ is an epimorphism of rings).
  

  According to (\ref{11/7/20.2}), $\bar{L}_{\shortparallel}$ is then associated with the special unital norm on $\bar{A}$ given by $\ell_{\rm sp}(\bar{s}) = [\text{det}(\underline{\mathcal{I}}\Omega\bar{s})]^{\frac{1}{m}}$ (we are abusing notation here, since $\ell_{\rm sp}(\bar{s})$ is obviously not the same function as $\ell_{\rm sp}(\underline{s})$ on the left-hand-side of (\ref{1/9/25.1}), as these special unital norms are associated with different algebras - but it is convenient to use the same notation for them, and no confusion arises).  That is, in evaluation of the integral resulting from  (\ref{11/7/20.2}), a path of integration from $\mathbf{1}_{\bar{A}}$ to $\bar{s}$ can be replaced by a path of integration from $\mathbf{1}_{\underline{A}}$ to $\underline{s}$ because of the projection implied by the integrand resulting from (\ref{1/9/25.2}).   Furthermore, on this path,
  \begin{eqnarray}\nonumber  [\bar{L}_{\shortparallel}\,\bar{t}^{-1}]\,\cdot\, \mathbf{d}\bar{t} &=& \left(\bar{t}^{-1}\right)^T \bar{L}_{\shortparallel}\,\mathbf{d}\bar{t} =  \frac{\|\mathbf{1}_{\bar{A}}\|^2}{\|\mathbf{1}_{\underline{A}}\|^2}\left(\bar{t}^{-1}\right)^T \Omega^T \underline{L}_{\shortparallel} \, \Omega \, \mathbf{d}\bar{t}  = \frac{\|\mathbf{1}_{\bar{A}}\|^2}{\|\mathbf{1}_{\underline{A}}\|^2}\,\left(\underline{t}^{-1}\right)^T\, \underline{L}_{\shortparallel}\,\mathbf{d}\underline{t} \\ &=& \frac{\|\mathbf{1}_{\bar{A}}\|^2}{\|\mathbf{1}_{\underline{A}}\|^2}\,[\underline{L}_{\shortparallel}\,\underline{t}^{-1}]\,\cdot\, \mathbf{d}\underline{t}.\label{2/5/25.2}\end{eqnarray}   In view of (\ref{2/5/25.2}), evaluation of (\ref{11/7/20.2}) then indicates that $\ell_{\rm sp}(\bar{s})$ is given by the right-hand-side of the first equality in (\ref{1/9/25.1}), and hence is given by the right-hand-side of the second equality in (\ref{1/9/25.1}). 
  

  It only remains to show that there always exists at least one such $\Omega$. 
  
  Evidently, either $\bar{A}$ is simple, or there exists an epimorphism $\Omega:\bar{A}\rightarrow \underline{A}$ where $ \underline{A}$ is a simple algebra.   By famous theorems of Wedderburn and Frobenius, a simple algebra is isomorphic to one of the real algebras $M_q(\mathbb{R})$, $M_q(\mathbb{C})$, or $M_q(\mathbb{H})$, for some positive integer $q$.  The left regular representation of each of these algebras is self-adjoint.  Thus, if $\bar{A}$ is simple, the theorem conclusion immediately follows by Lemma \ref{5/13/22.1}.  However, if $\bar{A}$ is not simple, the theorem conclusion follows from the first sentence of this paragraph and the first two paragraphs of this proof.
  \end{proof}

 Half of the examples of associative algebras $\bar{A}$ collected in Table III of \mbox{Section 7} are such that their left regular representations are not self-adjoint as subalgebras of $M_n(\mathbb{R})$, thereby illustrating Theorem \ref{11/17/22.1}.  Calculations related to such algebras are explicitly highlighted in the examples presented in Section 7.5 Section 7.6.3.


  \section{Combinatorial issues}
Regarding computation of the invariants identified in Corollary \ref{8/18/23.2}, we can make the following observations.  As indicated in the proof of Proposition \ref{9/2/24.1}, computation of $\mathfrak{u}$ in the form of $M_\mathfrak{u}$ is straightforward, as is the accompanying computation of the dimension of $\mathfrak{u}$.   Computation of the largest rank of a member of $\mathfrak{u}$ is also easy, e.g., using the method of minors.   Evaluation of the number of parameters to which ${\rm det}(M_{\mathfrak{u}})$ is sensitive is simply a matter of inspection.  Examples of invariants derived from $V_\mathbb{R}({\rm det}(M_{\mathfrak{u}}))$ include the real dimension of that real hypersurface (for which an efficient algorithm is \mbox{available \cite{bannwarth:2015})}, and the number of irreducible components of $V_\mathbb{R}({\rm det}(M_{\mathfrak{u}}))$.  

 
Based on the use of one or more simultaneously selected computable invariants identified from Corollary \ref{8/18/23.2}, various more-or-less combinatorial questions arise when the invariants (singly or jointly) are considered in combination with the invariant given by algebra dimension.  For example, with respect to a chosen algebra dimension,  


\begin{enumerate} \item How many values of the anti-rotor dimension occur?  \item How many irreducible components of the various $V_\mathbb{R}({\rm det}(M_{\mathfrak{u}_A}))$ can occur?  \item How many algebra isomorphism classes are there such that ${\rm det}(M_{\mathfrak{u}_A})$ is identically $0$?  \item For each $k\le \text{dim}(\mathfrak{u}_A)$, can $k$ ever be the maximal rank of a member of $\mathfrak{u}_A$?   \end{enumerate}  
 Quest for the above characterizations of the set of $n$-dimensional algebras can be considered a response to the intractability of the algebra classification problem.  On the other hand, it is likely that many of these combinatorial problems will be very difficult.

Concerning the general question of the distribution of $n$-dimensional algebras into distinct ``bins" (i.e., classes) defined by the anti-rotor dimension, according to the proof of Proposition \ref{9/2/24.1}, $\mathfrak{u}$ is derived from an initial set of $\frac{n(n+1)}{2}$ parameters, which are then subjected to  $\frac{n(n-1)}{2}$ constraint equations.  One might then expect to be finally left with $\frac{n(n+1)}{2} - \frac{n(n-1)}{2} = n$ parameters - i.e., $\text{dim}(\mathfrak{u}) = n$ (so there would only be a single nonempty bin).  Of course, this reasoning is faulty.  Consider the analogous problem of finding the symmetric matrices $L$ that satisfy (\ref{5/22/25.3}) except that $s^{-1}$ is replaced with $s$ in that equation.  One quickly discovers that every constraint equation reduces to $0=0$, so that {\it every} symmetric $(n\times n)$-matrix is a solution (the dimension of the solution space is $\frac{n(n+1)}{2}$).  So, if some of the constraint equations for uncurling metrics were subject to  that kind of thing, or if the $\frac{n(n-1)}{2}$ constraints did not constitute a set of $n$ independent constraints, perhaps $\text{dim}(\mathfrak{u})$ could be some number greater than $n$.  Conversely, it might occur that a constraint equation eliminates more than one parameter, so that $\text{dim}(\mathfrak{u})$ turns out to be less than $n$ (along these lines, see the proof of Proposition \ref{4/8/22.1}).  On the other hand, we still might guess that $\text{dim}(\mathfrak{u}) = n$ would define the most popular bin.  In fact, for all three isomorphism classes of two-dimensional algebras, Table I of Section 7 indicates that \mbox{$\text{dim}(\mathfrak{u}) = 2$}.  For the three-dimensional algebras, Table II indicates that five of the six algebra isomorphism classes satisfy $\text{dim}(\mathfrak{u}) = 3$, although the sixth satisfies $\text{dim}(\mathfrak{u}) = 2$.  Moreover, as regards the four-dimensional algebras, it turns out that $\text{dim}(\mathfrak{u}_{M_2(\mathbb{R})}) = 1$ (as follows from Proposition \ref{4/8/22.1}). So a precise understanding of the distribution of $n$-dimensional algebras into the bins/classes defined by $\text{dim}(\mathfrak{u})$ remains a goal.  

However, Corollary \ref{8/17/22.1} is a first general result relevant to a different aspect of the distribution of $n$-dimensional algebras into bins/classes defined by the anti-rotor dimension, as it indicates that if the anti-rotor is one-dimensional and has a nonsingular member then the algebra is simple.  

Another way to look at all this is as follows.
For a given integer $n\ge 4$, there are uncountably many isomorphism classes of $n$-dimensional algebras.  To deal with the accompanying unwieldy situation as regards classification, it becomes of interest to instead pursue combinatorial questions regarding the above discrete invariants as a means of characterizing the {\it set} of algebras as an entity in itself.   For example, consider the characterization defined by the lattice subset, \[ \mathcal{X}_{\text{dim}(\mathfrak{u})} \equiv \big\{\big(\text{dim}(A),\text{dim}(\mathfrak{u}_A)\big): A \text{ is an algebra}\big\},\] where $\text{dim}(\cdot)$ indicates dimension.  If computable, $\mathcal{X}_{\text{dim}(\mathfrak{u})}$ would be a depiction of whether the set of algebras with dimension $n$ and anti-rotor dimension $m$ is empty or not, for any $(n,m)$.  That is, algebra dimension and anti-rotor dimension are treated in combination, and counting is with respect to 0 (empty) and 1 (not empty). Although the possibility of a full determination of $\mathcal{X}_{\text{dim}(\mathfrak{u})}$ is in doubt, one can hope to reveal many of its distinguishing features (some perhaps in an  asymptotic or order of magnitude sense, as $n$ becomes large). 
We already know some simple things like, \begin{enumerate}\item For $h(n)$ equal to the cardinality of $\{m\in\mathbb{Z}:(n,m)\in \mathcal{X}_{\text{dim}(\mathfrak{u})}\}$, we have $h(n)\le  1+\frac{n(n+1)}{2}$ (which is obvious). \item $\mathcal{X}_{\text{dim}(\mathfrak{u})} \subset \mathbb{N}\times\mathbb{N}$ (Proposition \ref{6/11/25.1}), \item $(n^2,1) \in \mathcal{X}_{\text{dim}(\mathfrak{u})}$ (Proposition \ref{4/8/22.1}).
\end{enumerate}
As already indicated, a problem of particular interest would be to provide a much better understanding of $h(n)$ than that supplied by the first item in the above list.

There are evidently many combinatorial questions that can be proposed to characterize the set of algebras, stemming from the invariants deriving from the anti-rotor (treated singly or jointly).


\section{Basic functorial aspects and a couple of useful results}
 Theorem \ref{6/6/23.1} indicates that the association of algebras with anti-rotors is functorial.  As a basic example, let {\bf A} be the category each of whose objects is an algebra, where the morphisms are algebra isomorphisms, and let {\bf S} be the category each of whose objects is a finite-dimensional vector space of real symmetric $(n\times n)$-matrices, where the value of $n\ge 1$ depends on the object (i.e., the members of the object vector space consist of different linear combinations of the members of a finite set of linearly independent real symmetric $(n\times n)$-matrices for some fixed value of $n$).  For objects $S_1,S_2$ of {\bf S}, there is a morphism $k$ if there exists a nonsingular real matrix $k$ such that \mbox{$S_1  = k^T  S_2 k$} (using the same symbol to denote the morphism and the associated matrix).  Obviously, the identity morphisms exist and the morphisms are associative.    Definition \ref{7/21/23.1} specifies a mapping from algebras (the objects in {\bf A}) to anti-rotors (which are objects in {\bf S}), and (keeping in mind Theorem \ref{6/6/23.1}) the equation in the penultimate sentence indicates how a morphism in {\bf A} selects a morphism in {\bf S}.  It is thus easily verified that (\ref{5/23/23.1}) leads to a functor from {\bf A} to {\bf S}.  This immediately results in a class of algebra isomorphism invariants, some of which are listed in  \mbox{Corollary \ref{8/18/23.2}}.
Continuing in that vein, this program can similarly address the algebra epimorphism problem (determining whether two algebras are related by an epimorphism) using a simple modification of the above functor.  

Thus, let ${\bf E}$ be the category each of whose objects is an algebra (i.e., a real unital associative algebra of nonzero finite-dimension), and a morphism $k$ referable to a source $A$ and target $A'$ exists if there is an algebra epimorphism  $k:A\rightarrow A'$ (obviously the identity morphisms exist, and the morphisms are associative).    

We define category {\bf P} as follows. An object is a symmetric $(n\times n)$-matrix where $n$ is any chosen positive integer, and where each matrix entry is a linear combination of the components of an $m$-tuple of distinct real parameters, $m\ge 1$.  Furthermore, this parametrized matrix has the feature that if all of the $m$ parameters are replaced by 0 except for one parameter being replaced by 1, and one forms the collection of all such matrices except that a different one of the $m$ parameters is replaced by $1$ in each case, then the resulting $m$ real matrices are linearly independent.  Thus, for any object of {\bf P}, the implied set of real matrices resulting from all realizations of the associated $m$-tuple of parameters is an $m$-dimensional vector space (a realization of the $m$-tuple of parameters is the $m$-tuple resulting from a choice of real values for the $m$ parameters).  An object of {\bf P} is considered to be the same as the latter vector space, each of whose members is a real matrix, and such an object can be represented by a matrix whose entries are linear combinations of a set of real parameters.   There is a morphism pointing from object $P$ to object $P'$ if there exists a real matrix $w$ such that $ww^T$ is nonsingular and $w^TPw$ is a subspace of $P'$.  So we can say that the latter morphism results from $w$.   We could similarly notate the morphism as $w$, but it turns out that would cause confusion once we introduce the functor.  Therefore, we will instead suggestively notate the morphism implied by $w$ as the symbol ``\,$w^T[\cdot]w$\,".    The composition of morphism $\begin{CD}  P @>w^T[\cdot]w>> P'\end{CD}$ with morphism $\begin{CD}  P' @>w'^T[\cdot]w'>> P''\end{CD}$ is the morphism with source $P$ and target $P''$ resulting from $ww'$.  This is truly a morphism since it is easy to verify $(ww')^T P(ww')\subset P''$ and $(ww')(ww')^T$ is nonsingular.
Given the way we are notating the morphisms, \mbox{we have} \begin{equation}\label{2/22/25.1}(w'^T[\cdot]w')\circ (w^T[\cdot]w) = (ww')^T[\cdot](ww').\end{equation} It is easy to verify that composition of morphisms is associative. The identity morphisms obviously exist. Thus, {\bf P} is a category.

We define a mapping $F$ from category ${\bf E}$ to category {\bf P} as follows.  First, for any object $A$ of {\bf E}, $F[
A] \equiv \mathfrak{u}_A$.  Second, given a morphism $k_1$ from $A_1$ to $A_2$, if $k_1^T[\mathfrak{u}_{A_2}]k_1 \subset \mathfrak{u}_{A_1}$ and $k_1k_1^T$ is nonsingular, then the morphism $F(k_1)$ pointing from $F[A_2]$ to $F[A_1]$ is taken to be $\begin{CD}  \mathfrak{u}_{A_2} @>k_1^T[\cdot]k_1>> \mathfrak{u}_{A_1}\end{CD}$.  Otherwise, $F(k_1)$ is undefined.

\begin{theorem}\label{6/13/22.1}  $F: \mathbf{E}\rightarrow \mathbf{P}$ is a contravariant functor. 
 \end{theorem}
 \begin{proof}  The identity morphisms are obviously preserved.
We will next show show that the category $\mathbf{E}$ morphism $ \begin{CD}  A_1 @>k_1>> A_2\end{CD}$ implies the existence of category {\bf P} morphism $ \begin{CD}  F[A_2] @>F(k_1)>> F[A_1]\end{CD}$.  That is, we will show that $k_1^T[\mathfrak{u}_{A_2}]k_1 \subset \mathfrak{u}_{A_1}$ with $k_1k_1^T$ nonsingular.

 Taking $s_1$ to be notation for a variable in $A_1$, it follows from the epimorphism assumption that $(k_1s_1)^{-1} = k_1(s_1^{-1})$.  By definition, given the anti-rotor $F[A_2]$ associated with $A_2$, and any real matrix  $L_2 \in F[A_2]$, we have $d\left((s_2^{-1})^TL_2 (\mathbf{d}s_2)\right) = 0$ in some neighborhood of ${\bf 1}_{A_2}$, where $s_2$ is notation for a variable in $A_2$.  Since $k_1$ is an epimorphism, we are entitled to specify $s_2=k_1 s_1$ with respect to appropriate neighborhoods of $\mathbf{1}_{A_1}$ and $\mathbf{1}_{A_2}$.   It follows that \begin{eqnarray}\nonumber 0 = d\left((s_2^{-1})^TL_2 \mathbf{d}s_2\right) &=&  d\left(((k_1s_1)^{-1})^TL_2 \mathbf{d}(k_1s_1)\right) \\  &=& d\left((s_1^{-1})^T \big[k_1^T L_2k_1\big]\mathbf{d}s_1\right), \label{7/4/22.5}  \end{eqnarray}   using the equation in the first sentence of this paragraph (the epimorphism assumption), and the fact that a differential commutes with a linear transformation.  Equation (\ref{7/4/22.5}) is simply an instance of (\ref{5/22/25.3}), so it follows that symmetric $k_1^T L_2k_1$ is an uncurling metric associated with $A_1$.  
 Consequently, $k_1^TF[A_2]k_1 = k_1^T\mathfrak{u}_{A_2}k_1$ is a subspace of $F[A_1] = \mathfrak{u}_{A_1}$.  Furthermore, $k_1k_1^T$ is nonsingular since $k_1$ is an epimorphism.  Hence, $k_1^T[\cdot]k_1$ is a morphism pointing from $F[A_2]= \mathfrak{u}_{A_2}$ to $F[A_1]=\mathfrak{u}_{A_1}$.  That is, $F(k_1)=k_1^T[\cdot]k_1$ is well defined.

Finally, consider
      \[  \begin{CD}  
A_1 @>k_1>> A_2 @>k_2>> A_3\\
 @VVV   @VVV  @VVV\\
\mathfrak{u}_{A_1}   @<k_1^T[\cdot]k_1<<  \mathfrak{u}_{A_2}  @<k_2^T[\cdot]k_2<<  \mathfrak{u}_{A_3} \\ 
 \end{CD} \]
Based on the prior paragraph, morphism $F(k_1) = k_1^T[\cdot]k_1$ exists and morphism $F(k_2) = k_2^T[\cdot]k_2$ exists.  From observations made in the penultimate paragraph preceding this theorem's statement, the composition of these morphisms then exists.  According to the prior paragraph, $F(k_2\circ k_1)$ also exists.   Thus, from (\ref{2/22/25.1}) and the definition of $F$, we have \[ F(k_2\circ k_1) = (k_2k_1)^T[\cdot](k_2k_1) = (k_1^T[\cdot] k_1)\circ (k_2^T[\cdot]k_2)  = F(k_1)\circ F(k_2).\]   The theorem then follows.
\end{proof}

This leads to some useful results.  



\begin{corollary} \label{8/17/22.5} If $k_1:A_1\rightarrow A_2$ is an epimorphism, then ${\rm dim}(\mathfrak{u}_{A_2}) \le {\rm dim}(\mathfrak{u}_{A_1})$. \end{corollary}  

\begin{proof} By Theorem \ref{6/13/22.1} and the definition of a morphism in $\mathbf{P}$, we have \mbox{$k_1^T \mathfrak{u}_{A_2} k_1 \subset \mathfrak{u}_{A_1}$} with $k_1k_1^T$ nonsingular.  It follows that  $k_1k_1^T \mathfrak{u}_{A_2} k_1k_1^T \subset k_1\,\mathfrak{u}_{A_1}k_1^T$.  Since $k_1k_1^T$ is nonsingular, we then have 
\begin{equation}\label{2/23/25.1} \mathfrak{u}_{A_2} \subset (k_1k_1^T)^{-1}k_1\, \mathfrak{u}_{A_1} k_1^T(k_1k_1^T)^{-1}.\end{equation}  However, \begin{equation}\label{2/23/25.2}  \text{dim}\left((k_1k_1^T)^{-1}k_1 \,\mathfrak{u}_{A_1} k_1^T(k_1k_1^T)^{-1}\right) \le \text{dim}(\mathfrak{u}_{A_1}).\end{equation} Combining (\ref{2/23/25.1}) and (\ref{2/23/25.2}), the corollary follows.
   \end{proof}
   
  Corollary \ref{8/17/22.5} is useful in excluding the possibility that two algebras are related by an epimorphism (e.g., in excluding an epimorphism from an algebra to one of its subalgebras without having to identify ideals; see the twelfth and thirteenth rows of Table III in Section 7.6 for an example).


\begin{corollary}\label{8/17/22.1} If the anti-rotor of an algebra is one-dimensional and contains a nonsingular member, then the algebra is simple. \end{corollary}

\begin{proof} Suppose algebra $A_1$ is $j$-dimensional and has a non-trivial ideal, and its anti-rotor is one-dimensional and has a nonsingular member.  It follows that we can write $\mathfrak{u}_{A_1} = F[A_1] = \{\alpha L_1:\alpha\in \mathbb{R}\}$, where $L_1$ is a some real nonsingular $(j\times j)$-matrix.  Since $A_1$ has a non-trivial ideal, there exists an algebra epimorphism $k_1:A_1\rightarrow A_2$ where $A_2$ has dimension $q>0$, and $k_1$ has a nontrivial kernel.  Together, \mbox{Corollary \ref{8/17/22.5}} and the fact that the dimension of an algebra's anti-rotor is \mbox{at least 1} (Proposition \ref{6/11/25.1}), indicate that the dimension of the anti-rotor of $A_2$ is 1.   Thus we can write $\mathfrak{u}_{A_2} = F[A_2] = \{\alpha L_2:\alpha\in \mathbb{R}\}$, where $L_2$ is a nonzero real $(q\times q)$-matrix.  It follows that $k_1^T L_2 k_1$ also cannot be the zero matrix - since if it was then $k_1k_1^T L_2 k_1k_1^T$ would be the zero matrix, and applying $(k_1k_1^T)^{-1}$ to both sides of the prior expression would indicate that $L_2$ {\it is} the zero matrix.  

From (\ref{7/4/22.5}), we know that $ k_1^T L_2 k_1$ must be a member of $F[A_1]$, and from the prior paragraph we know that $ k_1^T L_2 k_1$ is not the zero matrix.  But the determinant of $k_1^T L_2 k_1$ is zero, since $k_1$ has a  nontrivial kernel.  And yet, the only member of the anti-rotor $F[A_1]$ having zero determinant is the zero matrix (corresponding to $\alpha=0$), since $L_1$ is nonsingular by assumption.   Thus, $k_1$ cannot exist, and so $A_1$ is simple. \end{proof}

Since according to Proposition \ref{4/8/22.1} the algebra of real $(n\times n)$-matrices $M_n(\mathbb{R})$ has a one-dimensional anti-rotor, a consequence of Corollary \ref{8/17/22.1} is an alternative demonstration that $M_n(\mathbb{R})$ is simple.  

Note that the converse of Corollary \ref{8/17/22.1} is not true.  For example, $\mathbb{C}$ is simple but has a two-dimensional anti-rotor (as noted in the first paragraph of Section 7.2). 
  
 \newpage
     
     \section{Illustrations} \subsection{Algebras with dimension less than 5.}
     
Taking $\sigma_1,\dots,\sigma_m$ to be the parameters appearing in $M_{\mathfrak{u}}$, the integral on the right-hand-side of (\ref{6/7/23.1}) can be written as 
\begin{equation}\label{9/26/23.1} \log  \mathfrak{u}^d = \int_{\mathbf{1}}^s [M_{\mathfrak{u}} (t^{-1})]\cdot \mathbf{d}t =\sum_{q=1}^m \left(\sigma_q \int_{\mathbf{1}}^s [\mathfrak{u}_q t^{-1}]\cdot \mathbf{d}t\right).\end{equation}
On the right-hand-side, $\mathfrak{u}_q$ is defined from $M_\mathfrak{u}$ by setting $\sigma_q =1$ with all other parameters set to zero.   All integrals are path-independent, and thus can be evaluated on a linear path of integration from \mbox{$\mathbf{1}$ to $s$} - on which the integrands are rational functions of one variable.  Thereby, the integration in principle yields a sum of rational function terms, logarithm terms, and arctangent terms.  Nevertheless, in dimensions greater than four a closed form expression for the indefinite integration (where $s$ is a variable) is not possible because it would depend on expressions for the roots of the polynomial in the denominator of the rational function integrand, which will depend on the indefinite components of $s$.  

However, in dimensions less than five the function obtained through evaluation of the integrals on the right-hand-side of (\ref{9/6/24.2}) is computable in closed form, and given by respective sums of rational function terms, logarithm terms, and arctangent terms each multiplied by one of the parameters $\sigma_q$ arising in $M_{\mathfrak{u}}$.  But the rational function terms can be collected into a {\it single} rational function term involving some number $\tau_{A,{\rm rat}}$ of distinct parameters.  Furthermore, the function's logarithm terms can be collected into a {\it single} (real) logarithm involving some number $\tau_{A,{\rm log}}$ of distinct parameters.  Similarly, its arctangent terms, which can be alternatively expressed in the format $\text{arctan}(w) = \frac{i}{2}\log\frac{1-iw}{1+iw}$, can be collected into a {\it single} such logarithm term utilizing imaginary coefficients, involving some number $\tau_{A,{\rm arc}}$ of distinct parameters. 
Given two isomorphic algebras $A_1$, $A_2$, Theorem \ref{6/6/23.2}  implies that based on dimensional considerations relating to the respective numbers of distinct parameters associated with the single rational function term, the single real logarithm term, and the single logarithm term with imaginary coefficients, we must have \begin{equation}\label{8/17/25.1} (\tau_{A_1,{\rm rat}},\tau_{A_1,{\rm log}}-1,\tau_{A_1,{\rm arc}}) = (\tau_{A_2,{\rm rat}},\tau_{A_2,{\rm log}}-1,\tau_{A_2,{\rm arc}}).\end{equation}
The above equation expresses three more algebra isomorphism invariants.
We have reduced the value of the second component of the ordered triplets by 1 since there is always at least one parameter multiplying a logarithm term, given by the 1-dimensional subspace of the anti-rotor's dual generated by multiples of a unital norm that is related to the usual norm as guaranteed by Theorem \ref{11/17/22.1}.  The equated  expressions in (\ref{8/17/25.1}) are referred to as ``reduced triples".
      

\subsection{Two-dimensional algebras}
Up to an isomorphism, there are only three two-dimensional algebras.
For the algebra of complex numbers, $\mathbb{C}$, we can write its elements $s$ as $(x,y)$, with the product given by $(x,y)(z,w) = (xz-yw, yz+xw)$, and $\mathbf{1}=(1,0)$, so that $\|\mathbf{1}\|^2=1$.  Also, $(x,y)^{-1} = (\frac{x}{x^2+y^2},\frac{-y}{x^2+y^2})$.  We start out by writing  $\mathfrak{u}_\mathbb{C} = \left[\begin{array}{ccc}
 \alpha & \beta   \\ 
 \beta & \gamma    
\end{array}\right]$.  \mbox{We then} obtain  $\mathfrak{u}_\mathbb{C}s^{-1} = \left[\begin{array}{ccc}
 \alpha & \beta   \\ 
 \beta & \gamma    
\end{array}\right]\left(\begin{array}{ccc}
x   \\ 
y   
\end{array}\right)^{-1} = \left(\begin{array}{ccc}
\frac{\alpha x - \beta y}{x^2+y^2}  \\ 
\frac{\beta x - \gamma y}{x^2+y^2} 
\end{array}\right)$.  As regards the right-hand-side of this equation, satisfaction of (\ref{5/22/25.3}) requires that the partial derivative with respect to $x$ of the second component must equal the partial derivative with respect to $y$ of the first component.  Consequently, $\mathfrak{u}_{\mathbb{C}} = \left[\begin{array}{ccc}
 \alpha & \beta   \\ 
 \beta & -\alpha.   
\end{array}\right]$.  Using Definition \ref{9/21/23.2}, a simple exercise in integration leads to $\log  \mathfrak{u}^d_{\mathbb{C}} = \frac{\alpha}{2}\log(x^2+y^2) + \beta{\rm arctan}(\frac{y}{x})$.   The invariant reduced triple (\ref{8/17/25.1}) is then (0,0,1).


The split-complex numbers $\slashed{\mathbb{C}}$ constitute an algebra for which the elements are  $(x,y)$ with product $(x,y)(z,w) = (xz+yw,xw+yz)$, $ \mathbf{1} = (1,0)$, and $\|\mathbf{1}\|^2 = 1$.  Also, $(x,y)^{-1} = (\frac{x}{x^2-y^2},\frac{-y}{x^2-y^2})$. \hspace{1cm}  We then find $\mathfrak{u}_{\slashed{\mathbb{C}}} = \left[\begin{array}{ccc}
 \alpha & \beta   \\ 
 \beta & \alpha.    
\end{array}\right]$, \hspace{.5cm} and \hfill \\ $\log  \mathfrak{u}^d_{\slashed{\mathbb{C}}} = \frac{\alpha}{2}\log[(x+
y)(x-y)] +  \frac{\beta}{2}\log \frac{x+y}{x-y}$.  In fact, this can be simplified to $\log  \mathfrak{u}^d_{\slashed{\mathbb{C}}} = \frac{\alpha+\beta}{2}\log(x+y) + \frac{\alpha-\beta}{2}\log(x-y)$.  Of course, $\frac{\alpha+\beta}{2}$, $\frac{\alpha-\beta}{2}$ are themselves two independent parameters that we could name $\alpha'$, $\beta'$.  Either way, the reduced triple (\ref{8/17/25.1}) is (0,1,0).

The final real two-dimensional unital algebra (up to an isomorphism) is the algebra of dual numbers $\mathbb{D}$.  The elements are $(x,y)$ with product given by $(x,y)(z,w) = (xz,xw+yz)$, with $\mathbf{1}=(1,0)$ and $\|\mathbf{1}\|^2 = 1$. Also, $(x,y)^{-1} = (\frac{1}{x},-\frac{y}{x^2})$.  We then find $\mathfrak{u}_\mathbb{D} = \left[\begin{array}{ccc}
 \alpha & \beta   \\ 
 \beta & 0    
\end{array}\right]$ and $\log  \mathfrak{u}^d_{\mathbb{D}} = \alpha\log x + \beta\left(\frac{y}{x}\right)$.  The reduced triple (\ref{8/17/25.1}) is (1,0,0).

   Now we consider $\mathbb{R}\times\mathbb{R}$ (understood to be a direct product of rings), which is isomorphic to $\slashed{\mathbb{C}}$.  We have elements as $(x,y)$, the product $(x,y)(z,w) = (xz,yw)$, and the multiplicative identity $\mathbf{1}=(1,1)$, so that $\|\mathbf{1}\|^2 = 2$.  Also $(x,y)^{-1} = (\frac{1}{x},\frac{1}{y})$.  We then find $\mathfrak{u}_{\mathbb{R}\times\mathbb{R}} = \left[\begin{array}{ccc}
 \alpha & 0   \\ 
 0 & \gamma    
\end{array}\right]$ and $\log  \mathfrak{u}^d_{\mathbb{R}\times\mathbb{R}} = \frac{\alpha}{2}\log x + \frac{\gamma}{2}\log y$.  The reduced triple (\ref{8/17/25.1}) is then (0,1,0).

$\frac{1}{2}\left[\begin{array}{ccc}
 1 & -1   \\ 
 1 & 1    
\end{array}\right]$ provides the isomorphism between the algebras $\slashed{\mathbb{C}}$ and $\mathbb{R} \times\mathbb{R}$.  It is easy to verify that their anti-rotors as identified above are consistent as regards to Theorem \ref{6/6/23.1}, and we have seen that the two algebras have the same reduced triple.  


Incidentally, if we multiply $\mathfrak{u}_{\mathbb{C}}$ on the left by the exchange matrix $\left[\begin{array}{ccc}
 0 & 1   \\ 
 1 & 0    
\end{array}\right]$, and also multiply $\mathfrak{u}_{\mathbb{D}}$ on the left by $\left[\begin{array}{ccc}
 0 & 1   \\ 
 1 & 0    
\end{array}\right]$, we obtain the left regular representations of $\mathbb{C}$ and $\mathbb{D}$, respectively.   In fact, $\mathfrak{u}_{\slashed{\mathbb{C}}}$ and $\mathfrak{u}_{\mathbb{R}\times\mathbb{R}}$ are already the same as the left regular representations of $\slashed{\mathbb{C}}$ and $\mathbb{R}\times\mathbb{R}$, respectively.  This type of phenomenon is {\it not} typical in higher dimensions, apart from special cases such as the algebras given by a direct product $\prod_{i=1}^k \mathbb{R}$, and the upper triangular Toeplitz matrix algebras.

 The three isomorphism classes of the two-dimensional algebras represented by $\mathbb{C}$, $\slashed{\mathbb{C}}$, and $\mathbb{D}$, respectively correspond to the reduced triples  $(0,0,1)$, $(0,1,0)$, and $(1,0,0)$ (i.e., pure elemental forms of the reduced triples associated with algebras in general).  There is only a single isomorphism class for the one-dimensional algebras and its reduced triple is (0,0,0).
 
Table I lists the three isomorphism classes of two-dimensional algebras,, along with their left regular representations, anti-rotors, and anti-rotor duals. 


\pagebreak
\hspace{0pt}
\vfill


  \centerline{\large The Three Isomorphism Classes of the Two-dimensional Algebras } 

\vspace{.5cm}

\hspace{.7cm}$ \begin{array}{cccccccc} 
\mbox{\underline{algebra}} & \mbox{\underline{anti-rotor $(\mathfrak{u})$}} & \mbox{\underline{unital norm group ($\mathfrak{u}^d$)}}  \\ 
\\
\mathbb{C} \approx \left[\begin{array}{ccc}
x  & -y    \\
 y & x \\   
\end{array}\right]   & \left[\begin{array}{ccc}
\alpha  & \beta    \\
 \beta & -\alpha \\   
\end{array}\right]  & (x^2+y^2)^{\frac{\alpha}{2}} \,e^{\beta\,{\rm arctan}\left(\frac{y}{x}\right)}  \\ \\
\slashed{\mathbb{C}} \approx \left[\begin{array}{ccc}
x  & y    \\
 y & x \\   
\end{array}\right]   & \left[\begin{array}{ccc}
\alpha  & \beta    \\
 \beta & \alpha \\   
\end{array}\right]  & [(x+y)(x-y)]^{\frac{\alpha}{2}} \left[\frac{x+y}{x-y}\right]^{\frac{\beta}{2}}  \\ \\ 
\mathbb{R}\times\mathbb{R} \approx \left[\begin{array}{ccc}
x  & 0    \\
 0 & y \\   
\end{array}\right]   & \left[\begin{array}{ccc}
\alpha  & 0    \\
 0 & \beta \\   
\end{array}\right]  & x^\frac{\alpha}{2} y^\frac{\beta}{2}  \\ \\
\mathbb{D} \approx \left[\begin{array}{ccc}
x  & 0    \\
 y & x \\   
\end{array}\right]   & \left[\begin{array}{ccc}
\alpha  & \beta    \\
 \beta & 0 \\   
\end{array}\right]  & x^\alpha\,e^{\beta \frac{y}{x}}  \\ \\
\end{array} $

\bigskip


\noindent {\bf Table I.} {\small In this table, as well as in Table II and Table III, an expression with Greek letters is understood to imply the set whose members are all realizations of those parameters as real numbers.  Each row indicates a representative  of one of the three two-dimensional algebra isomorphism classes along with its left regular representation, followed by the associated anti-rotor and unital norm group (``$\approx$" means \mbox{``isomorphic to"}). Note that  $\slashed{\mathbb{C}}$ and $\mathbb{R}\times\mathbb{R}$ are isomorphic, their respective anti-rotors are compatible as regards \mbox{Theorem \ref{6/6/23.1}} (via the algebra isomorphism $\frac{1}{2}\left[\begin{array}{ccc}
 1 & 1   \\ 
 -1 & 1    
\end{array}\right]$), and their unital norm groups are  compatible as regards the reduced triples (i.e., $(x+y)^{\frac{\alpha+\beta}{2}}(x-y)^{\frac{\alpha-\beta}{2}}$ versus $x^\frac{\alpha}{2} y^\frac{\beta}{2}$).}

\vspace{1.75cm} 

\subsection{Three-dimensional algebras.} The six isomorphism classes of  algebras having three dimensions are also instructive.  The structure constants for a representative of each of these are presented in \cite{kobayashi:2021}.  Table II lists those six algebras (in the order given in \cite{kobayashi:2021}) in matrix representations along with their respective anti-rotors and unital norm groups. The  reduced triples are, respectively, (0,2,0), (0,1,1), (1,1,0), (2,0,0), (2,0,0), and (0,1,0).



\hspace{0pt}
\vfill
\pagebreak
\vfill
\hspace{0pt}

  \centerline{\large The Six Isomorphism Classes of the Three-dimensional Algebras }

\vspace{.5cm}


\hspace{-2.4cm}$ \begin{array}{cccccccc} 
\mbox{\underline{algebra}} & \mbox{\underline{anti-rotor ($\mathfrak{u}$) }} & \mbox{\underline{unital norm group ($\mathfrak{u}^d$)}}  \\ 
\\
 \mathbb{R}\times \mathbb{R}\times \mathbb{R}  \approx \left[\begin{array}{ccc}
x  & 0 & 0   \\
 0 & y & 0 \\
 0 & 0 & z    
\end{array}\right]  \approx \mathbb{R}\times\slashed{\mathbb{C}}  &  \mathbb{R}\oplus \mathfrak{u}_{\mathbb{R}\times\mathbb{R}}, \begin{array}{ccc}
V_\mathbb{R}(\alpha\beta\gamma)   \\ 
(3,3,1,3,2,3)   
\end{array} & \begin{array}{ccc} x^{\frac{\alpha}{3}}y^{\frac{\beta}{3}}z^{\frac{\gamma}{3}} \\ (0,2,0) \end{array}
 \\ \\
   \mathbb{R}\times\mathbb{C} \approx \left[\begin{array}{ccc}
x  & 0 & 0   \\
 0 & y & -z \\
 0 & z & y    
\end{array}\right]   \hfill &  \mathbb{R}\oplus \mathfrak{u}_{\mathbb{C}}, \begin{array}{ccc}
V_\mathbb{R}(-\alpha(\beta^2+\gamma^2))   \\ 
(3,3,1,3,2,2)   
\end{array} & \begin{array}{ccc}  x^\alpha (y^2+z^2)^{\frac{\beta}{2}} \,e^{\gamma\,{\rm arctan}\left(\frac{z}{y}\right)} \\ (0,1,1)\end{array} \\ \\
   \mathbb{R}\times\mathbb{D} \approx \left[\begin{array}{ccc}
x  & 0 & 0   \\
 0 & y & 0 \\
 0 & z & y    
\end{array}\right]  \hfill &  \mathbb{R}\oplus \mathfrak{u}_{\mathbb{D}}, \begin{array}{ccc}
V_\mathbb{R}(-\alpha\gamma^2)   \\ 
(3,3,1,2,1,2)   
\end{array}  &  \begin{array}{ccc} x^\alpha y^{\beta} e^{\gamma\frac{y}{z}} \\ (1,1,0)\end{array} \\ \\
   \mathbb{T}_3  = \left[\begin{array}{ccc}
x  & y & z   \\
 0 & x & y \\
 0 & 0 & x    
\end{array}\right] \approx \left[\begin{array}{ccc}
x  & 0 & 0   \\
 y & x & 0 \\
 z & y & x    
\end{array}\right] \hfill &  \left[\begin{array}{ccc}
\alpha  & \beta & \gamma   \\
 \beta & \gamma & 0 \\
 \gamma & 0 & 0   
\end{array}\right], \begin{array}{ccc}
V_\mathbb{R}(\gamma^3)   \\ 
(3,3,1,1,0,1)   
\end{array}  & \begin{array}{ccc} x^\alpha e^{ \beta\frac{y}{x}  +\gamma \frac{z}{x} -\gamma \frac{ 1}{2}\left(\frac{y}{x}\right)^2 } \\ (2,0,0)\end{array} \\ \\
\left[\begin{array}{ccc}
x  & y & 0   \\
 0 & x & 0 \\
 0 & z & x    
\end{array}\right] \approx \left[\begin{array}{ccc}
x  & 0 & 0   \\
 y & x & 0 \\
 z & 0 & x    
\end{array}\right] \hfill &    \left[\begin{array}{ccc}
\alpha  & \beta & \gamma   \\
 \beta & 0 & 0 \\
 \gamma & 0 & 0   
\end{array}\right], \begin{array}{ccc}
V_\mathbb{R}(0)   \\ 
(3,2,1,0,3,0)   
\end{array}   &  \begin{array}{ccc} x^\alpha e^{ \beta\frac{y}{x}  +\gamma \frac{z}{x}  } \\ (2,0,0)\end{array} \\ \\
\left[\begin{array}{ccc}
x  & z & y   \\
 0 & x+y & 0 \\
 y & -z & x    
\end{array}\right]  \approx \left[\begin{array}{ccc}
x  & y & 0   \\
 y & x & 0 \\
 z & z & x-y    
\end{array}\right]  \hfill & \hspace{.5cm} \left[\begin{array}{ccc}
\alpha  & \beta & 0   \\
 \beta & \alpha & 0 \\
 
 0 & 0 & 0   
\end{array}\right] = \mathfrak{u}_{\slashed{\mathbb{C}}}\oplus \{0\}, \begin{array}{ccc}
V_\mathbb{R}(0)   \\ 
(2,2,1,0,2,0)
\end{array}  &  \begin{array}{ccc}(x+y)^\alpha (x-y)^\beta \\ (0,1,0)\end{array} \\ \\
\end{array} $



\vspace{.5cm}


\noindent {\bf Table II.} {\small   
Each row lists a representative  of one of the six three-dimensional algebra isomorphism classes along with its left regular representation as a subalgebra of $M_3(\mathbb{R})$, followed by the associated anti-rotor and unital norm group. The first matrix listed in \mbox{row four}, row five, and row six, is a matrix that indicates the manner in which the algebra elements multiply, and the second matrix is the left regular representation of that three dimensional algebra.  $\mathbb{T}_3$ is the three-dimensional upper triangular Toeplitz matrix algebra. The invariant ordered sextuple in each row refers to the invariants listed in Corollary \ref{8/18/23.2}, listing in order the the value of the dimension of $\mathfrak{u}$, the largest rank of a member of $\mathfrak{u}$, the smallest rank of a nonzero member of $\mathfrak{u}$, the dimension of the parameter subspace to which det$(M_{\mathfrak{u}})$ is sensitive, the real dimension of the real algebraic set defined by det$(M_{\mathfrak{u}}) =0$, and the number of irreducible components of $V_\mathbb{R}(\text{det}(M_{\mathfrak{u}}))$ (this hypersurface is indicated above the sextuple).  The reduced triple is indicated below each unital norm group.  In contrast to the two-dimensional algebras, not all anti-rotors of the three-dimensional algebras have the same dimension as the algebra.}
 

\vfill
\hspace{0pt}
\pagebreak



 \subsection{The algebra with vector space of elements $\mathbb{R}^n$ and component-wise multiplication} A higher dimensional analogue of the direct product of rings \mbox{$\mathbb{R}\times \mathbb{R}$}, i.e., $\prod_{i=1}^n \mathbb{R}$, is associated with an $(n-1)$-parameter family of  special unital norms.

 This algebra is isomorphic to the subalgebra of $ M_n(\mathbb{R})$  whose members are diagonal matrices. In this example, ${\bf 1} = (1,\dots,1)$ and $\|{\bf 1}\|^2 = n$.   It is easily appreciated that (\ref{5/22/25.3}) forwards the anti-rotor ${\rm diag} \{ \sigma_1, \sigma_2,\dots,\sigma_n\}$, so the anti-rotor is the same as the members of the left regular representation of the algebra.  This means that for $s=(x_1,\dots,x_n)$, we must have $L[(x_1,x_2,\dots,x_n)^{-1}] = (\frac{\sigma_1}{x_1}, \frac{\sigma_2}{x_2},\dots,\frac{\sigma_n}{x_n})$.     Evaluation of (\ref{11/7/20.2}) indicates that the members of the unital norm group are given by  \begin{equation}\label{1/15/22.5} \ell_{\sigma_1,\dots,\sigma_n}(x_1,\dots,x_n) = (x_1^{\sigma_1} x_2^{\sigma_2} \dots x_n^{\sigma_n})^\frac{1}{n}. \end{equation} From (\ref{7/7/21.1}), the special unital norms are given by (\ref{1/15/22.5}) under the constraint $\sum_1^n \sigma_i = n$. Thus, for this algebra we have a $(n-1)$-dimensional special unital norm family.   
 
 The geometric mean is the special unital norm given by (\ref{1/15/22.5}) with $\sigma_1=\sigma_2=\dots =\sigma_n = 1$.  
   
 \subsection{Upper triangular matrix algebras with independent entries on the main diagonal} Consider the real three-dimensional algebra with elements $(x,y,z)$ which multiply according to the matrices 
$\left[
\begin{array}{ccc}
 x & z   \\
 0 & y    
\end{array}
\right]$.  Or in other words, the algebra is $\left\{ \left[
\begin{array}{ccc}
 x & z   \\
 0 & y    
\end{array}
\right],x,y,z\in\mathbb{R}\right\}$.  We have $(x,y,z)^{-1} = \left( \frac{1}{x}, \frac{1}{y}, -\frac{z}{xy} \right)$, and ${\bf 1}= (1,1,0)$, and $\|{\bf 1}\|^2 = 2$.  Starting with a generic symmetric $L = \left[
\begin{array}{ccc}
 \alpha & \beta & \gamma   \\
 \beta & \delta & \epsilon \\ \gamma & \epsilon & \xi    
\end{array}
\right]$, 
it is required that the curl of $L(x,y,z)^{-1}$, which is \[\left( \frac{\epsilon}{xy} - \frac{\epsilon}{y^2} + \frac{\xi z}{xy^2},\,\,\, \frac{\gamma}{x^2} - \frac{\gamma}{xy} - \frac{\xi z}{x^2y}, \,\,\, \frac{\beta}{y^2} - \frac{\beta}{x^2} - \frac{\gamma z}{xy^2} + \frac{\epsilon z}{x^2y} \right),\] be such that each of its components vanishes in a neighborhood of $(1,1,0)$.  This implies that all entries of the matrix $L$ must be zero except for $\alpha$ and $\delta$.   Thus, $L(x,y,z)^{-1} = (\frac{\alpha}{x}, \frac{\delta}{y},0)$.  From (\ref{7/7/21.1}), we have the requirement that a normalized uncurling metric must satisfy \[ (x,y,z)\cdot [L(x,y,z)^{-1}] = \alpha +\delta =  2.\]  Equivalently, we can alternatively specify $\alpha \equiv 1+\sigma$ and $\delta \equiv 1-\sigma$, for arbitrary $\sigma\in\mathbb{R}$.   We thus have the normalized uncurling metric family composed of all  \begin{equation}\label{11/17/22.2} L = \left[
\begin{array}{ccc}
 1+\sigma & 0 & 0   \\
 0 & 1-\sigma & 0 \\ 0 & 0 & 0    
\end{array}
\right],\end{equation} $\sigma\in\mathbb{R}$.  In this case, each member of the family is singular.  Equation (\ref{11/7/20.2}) then leads to the 1-parameter family of special unital norms, \begin{equation}\label{12/21/21.1} \ell_\sigma(x,y,z) = x^\frac{1+\sigma}{2} y^\frac{1-\sigma}{2} = (xy)^{\frac{1}{2}}\left(\frac{x}{y}\right)^{\frac{\sigma}{2}}.\end{equation} 

The left regular representation of an element $(x,y,z)$ is given by $\left[\begin{array}{ccc}
x & 0 & 0   \\
 0 & y & 0 \\ 0 & z & x    
\end{array} \right]$.  Thus, the $M_3(\mathbb{R})$-transpose operation maps outside the algebra's representation.  However, there is an epimorphism to the two-dimensional subalgebra defined by elements $(x,y,0)$, and the latter's left regular representation is 
 $\left[\begin{array}{ccc}
x & 0  \\
 0 & y  
\end{array} \right]$.  The usual norm of the subalgebra is thus $xy$, and one of the special unital norms on the subalgebra is $(xy)^{\frac{1}{2}}$, and these differ only by an exponent - consistent with Lemma \ref{5/13/22.1}.   Note that the latter unital norm on the subalgebra provides the values for a special unital norm on the full algebra, corresponding to the choice of normalized uncurling metric in (\ref{11/17/22.2}) given by $\sigma = 0$, i.e., $\ell_0(x,y,z) = (xy)^{\frac{1}{2}}$, consistent with \mbox{Theorem \ref{11/17/22.1}}.  The fact that the transpose operation on $M_3(\mathbb{R})$ does not map the above left regular representation of $\left\{\left[\begin{array}{ccc}
x & z  \\
 0 & y  
\end{array} \right]:(x,y,z)\in \mathbb{R}^3\right\}$ to itself explains why there is no unital norm that is a power of the usual norm $x^2y$.
  
\subsection{Some algebras with unital norms that are sensitive to more element components than is the usual norm}

The usual norm on $ M_n(\mathbb{R})$ is sensitive to all components of an element of the algebra (because the determinant of an element of $M_n(\mathbb{R})$  is sensitive to the values of all entries of its matrix).    However, this is not true in general for associative algebras.  For example, for any algebra isomorphic to a triangular matrix algebra, the usual norm is only sensitive to the components on the main diagonal of the triangular matrix representation.  This is also true for the unital norm if all entries on the main diagonal of the triangular matrix representation are allowed to be distinct.  But if  at least two entries on the main diagonal always take equivalent values for the particular algebra,  then almost all of its unital norms family members will be sensitive to one or more components not on the main diagonal, in addition to all of the components on the main diagonal (these are matrix algebras for which an element that is not a diagonal matrix is never diagonalizable).  

It is easily seen from the presentation in Section 7.2 that the dual numbers are a two-dimensional example of this, and it
generalizes to higher dimensions as presented in Appendix A.  The next three examples highlight this feature.  

\subsubsection{The three-dimensional Toeplitz matrix algebra, $\mathbb{T}_3$} Consider the real algebra with elements  $(x,y,z)$ that multiply according to the matrices $ \left[
\begin{array}{ccc}
x & y & z   \\
0 & x & y \\
0 & 0 & x  
\end{array}
\right]$.    The constraints (\ref{7/21/23.1}) and (\ref{7/7/21.1}) lead to the normalized uncurling metrics, $ L = \left[
\begin{array}{ccc}
1 & \beta & \gamma   \\
\beta & \gamma & 0 \\
\gamma & 0 & 0  
\end{array}
\right]$, $\beta,\gamma\in\mathbb{R}$, which is just the $n=3$ case of Corollary \ref{7/1/22.1}.    Evaluating (\ref{11/7/20.2}), we now have a 2-parameter family of special unital norms, \begin{equation}\label{3/31/21.3} \ell_{\beta,\gamma}(x,y,z) = x\, e^{ \beta\frac{y}{x}  +\gamma \frac{z}{x} -\gamma \frac{ 1}{2}\left(\frac{y}{x}\right)^2 }.\end{equation}    The choice $\beta = \gamma = 0$ produces the special unital norm, $\ell_{0,0}(x,y,z)= x$, which is comparable to the usual norm, $x^3$.  
 The other special unital norms are distinguished by their dependence on $y$ and/or $z$, i.e., unlike the usual norm.  The non-units are associated with a different type of unital norm singularity (as $x$ approaches 0) than seen in examples considered so far other than $\mathbb{D}$.
 

 \subsubsection{A four-dimensional algebra} Consider the real unital algebra with elements 
$(x,y,z,w)$ which multiply according to the matrices 
$\left[
\begin{array}{ccccc}
x & y & w   \\
0 & x & z \\
0 & 0 & x  
\end{array}
\right]$.
Straightforward application of the constraints  (\ref{7/21/23.1}) and (\ref{7/7/21.1}) leads to the normalized uncurling metric family 
$\left[
\begin{array}{cccc}
1 & \beta & \gamma & \delta   \\
\beta & 0 & \frac{1}{2}\delta & 0 \\
\gamma & \frac{1}{2}\delta &  0 & 0   \\
\delta & 0 &  0 & 0
\end{array}
\right]$, and the associated special unital norm family,
\[ \ell_{\beta,\gamma,\delta}(x,y,z,w) = x\, e^{\beta\frac{y}{x} + \gamma\frac{z}{x} +\delta\frac{w}{x}- \frac{\delta}{2}\left(\frac{yz}{x^2}\right)}.\]


\subsubsection{A five-dimensional algebra}Consider the real unital algebra with elements 
$(x,y,z,v,w)$ which multiply according to the matrices
$\left[
\begin{array}{ccccc}
x & z & w   \\
0 & x & v \\
0 & 0 & y  
\end{array}
\right]$.
Straightforward application of the constraints (\ref{7/21/23.1}) and (\ref{7/7/21.1}) leads to the normalized uncurling metric family 
$\left[
\begin{array}{ccccc}
1+\sigma & 0 & \beta & 0  & 0   \\
0 & 1-\sigma & 0 & 0 & 0 \\
\beta & 0 &  0 & 0 & 0  \\
 0 & 0 &  0 & 0 & 0  \\
0 & 0 &  0 & 0 & 0
\end{array}
\right]$, and the associated special unital norm family 
\[ \ell_{\sigma,\beta}(x,y,z,v,w) =  (xy)^{\frac{1}{2}}\left(\frac{x}{y}\right)^{\frac{\sigma}{2}} e^{\frac{\beta}{2} \frac{z}{x}}.\]

In this example, almost all unital norms are influenced the element's $z$-component, while this sensitivity is not present with the usual norm (given by $x^4y$).  However, like the usual norm,  the unital norm remains insensitive to the $w$ and $v$ components.  Once again, there is a new unital norm singularity type as non-units are approached.  

Finally, the left regular representation of this algebra with elements $(x,y,z,v,w)$ is given by $ \left[
\begin{array}{ccccc}
x & 0 & 0 & 0 & 0   \\
0 & y & 0 & 0 & 0 \\
z & 0 &  x & 0 & 0  \\
v & 0 &  0 & y & 0  \\
w & 0 &  v & 0 & y
\end{array}
\right].$ 
The transpose operation on $M_5(\mathbb{R})$ does not map this representation algebra into itself, due to its effect on the last three columns and last three rows - which explains why there is no unital norm that is a power of the usual norm $x^2y^3$.  However, there is an epimorphism to the two-dimensional subalgebra with elements $(x,y,0,0,0)$.  The subalgebra's left regular representation is $\left[\begin{array}{ccc}
x & 0  \\
 0 & y  
\end{array} \right]$, which leads to the special unital norm $(xy)^\frac{1}{2}$ according to \mbox{Theorem \ref{11/17/22.1}}.


\hspace{-3.5cm} $ 
\begin{array}{cccccccc}
{\rm \underline{algebra}} & \mbox{\underline{usual norm}} & \mbox{\underline{special unital norm family}} & \mbox{\underline{anti-rotor ($\mathfrak{u}$) }} &  \bigskip \\ 
M_n(\mathbb{R}) & ({\rm det}(s))^n & ({\rm det}(s))^{\frac{1}{n}} & \alpha \,(n^2\times n^2)\mbox{\rm -``transpose\, matrix"}  \\  \\ 
\mathbb{C} & x^2+y^2 = r^2 & \sqrt{x^2+y^2}\,e^{\beta{\rm arctan}\left(\frac{y}{x}\right)} = re^{\beta\theta} & \left[\begin{array}{ccc}
 \alpha & \beta   \\ 
 \beta & -\alpha    
\end{array}\right]  =  \slashed{\mathbf{1}}_2 \, \mathbb{C}_\text{left} \\ \\ \slashed{\mathbb{C}}
 & x^2-y^2 & \begin{array}{cc}
 \sqrt{x^2-y^2}\, e^{\beta{\rm arctanh}\left(\frac{y}{x}\right)}   \\
 = (x^2-y^2)^{\frac{1}{2}}\left( \frac{x+y}{x-y}\right)^\frac{\beta}{2}    
\end{array} & \left[\begin{array}{ccc}
 \alpha & \beta   \\
 \beta & \alpha    
\end{array}\right]  = \slashed{\mathbb{C}}_\text{left} \\ \\
\mathbb{R}\times\mathbb{R} & xy & (xy)^{\frac{1}{2}}\left(\frac{x}{y}\right)^{\frac{\sigma}{2}} & \left[\begin{array}{ccc}
\sigma_1 & 0   \\
 0 & \sigma_2    
\end{array}\right] = [\mathbb{R}\times\mathbb{R}]_\text{left}  \\  \\
\vspace{-1mm} \mathbb{D}  & x^2 & x\, e^{\beta\frac{y}{x}} & \left[\begin{array}{ccc}
 \alpha & \beta   \\
 \beta & 0    
\end{array}\right] = \slashed{\mathbf{1}}_2\mathbb{D}_\text{left} \\ \\

\vspace{-1mm} \mathbb{H}  & (x^2+y^2+z^2+w^2)^2 & \sqrt{x^2+y^2+z^2+w^2} &  \,{\rm diag}\{\alpha,-\alpha,-\alpha,-\alpha\}   \\ \\

\mbox{\rm Cayley-Dixon seq., } \mbox{after }\mathbb{H} & (\mbox{\rm Euclidean norm})^2 & {\rm Euclidean\, Norm} &  \,{\rm diag}\{\alpha,-\alpha,-\alpha,\dots,-\alpha\} &  \\ \\
\mbox{Formally real Spin Factor Jordan} & (\mbox{\rm Minkowski norm})^2 & \mbox{\rm Minkowski norm} &  \,{\rm diag}\{\alpha,\alpha,\dots,\alpha\} & \\ \\ \prod_{i=1}^n \mathbb{R}
& \prod x_i & \left(\prod x_i^{\sigma_i}\right)^\frac{1}{n}, \sum\sigma_i=n & {\rm diag} \{\sigma_1,\dots,\sigma_n\} & \\ \\
\left\{ \left[\begin{array}{ccc}
 x & z   \\
 0 & y    
\end{array}\right]: (x,y,z)\in \mathbb{R}^3 \right\} & x^2y & (xy)^{\frac{1}{2}}\left(\frac{x}{y}\right)^{\frac{\sigma}{2}} & \left[\begin{array}{ccc}
\sigma_1  & 0 & 0   \\
 0 & \sigma_2 & 0 \\
 0 & 0 & 0    
\end{array}\right] & \\ \\
\left\{  \left[\begin{array}{ccc}
x  & y & z   \\
 0 & x & y \\
 0 & 0 & x    
\end{array}\right]: (x,y,z)\in \mathbb{R}^3   \right\}& x^3 & x\, e^{ \beta\frac{y}{x}  +\gamma \frac{z}{x} -\gamma \frac{ 1}{2}\left(\frac{y}{x}\right)^2 } &  \left[\begin{array}{ccc}
\alpha  & \beta & \gamma   \\
 \beta & \gamma & 0 \\
 \gamma & 0 & 0    
\end{array}\right]  \\ \\ \left\{ \left[
\begin{array}{ccccc}
x & y & w   \\
0 & x & z \\
0 & 0 & x  
\end{array}
\right] :  (x,y,z,w)\in\mathbb{R}^4 \right\} & x^4 & x\, e^{\beta\frac{y}{x} + \gamma\frac{z}{x} +\delta\frac{w}{x}- \frac{\delta}{2}\left(\frac{yz}{x^2}\right)}
& \left[
\begin{array}{cccc}
\alpha & \beta & \gamma & \delta   \\
\beta & 0 & \frac{1}{2}\delta & 0 \\
\gamma & \frac{1}{2}\delta &  0 & 0   \\
\delta & 0 &  0 & 0
\end{array}
\right]\\ \\
\left\{  \left[\begin{array}{ccc}
x  & z & w   \\
 0 & x & v \\
 0 & 0 & y    
\end{array}\right]:  (x,y,z,v,w)\in\mathbb{R}^5 \right\}  & x^4y & (xy)^{\frac{1}{2}}\left(\frac{x}{y}\right)^{\frac{\sigma}{2}}  e^{\frac{\beta}{2}  \frac{z}{x}}  & \left[\begin{array}{ccccc}
\sigma_1  & 0 & \beta & 0 & 0   \\
 0  & \sigma_2 & 0 & 0 & 0  \\
\beta   & 0 & 0 & 0 & 0 \\
 0  & 0 & 0 & 0 & 0  \\
 0  & 0 & 0 & 0 & 0    
\end{array}\right]  \\  \mathbb{T}_n & x_1^n & x_1 e^{P_{\gamma_2,\dots,\gamma_n}\left(\frac{x_2}{x_1},\dots,\frac{x_n}{x_1}\right)} & \slashed{\mathbf{1}}_n[\mathbb{T}_n]_\text{left} & \bigskip \end{array} 
$

\bigskip

\enlargethispage{1.5cm}
\noindent {\bf Table III.} {\small Miscellaneous examples. $\prod_{i=1}^n \mathbb{R}$ is the direct product of rings.  $\mathbb{T}_n$ is the $n$-dimensional upper triangular Toeplitz matrix algebra (treated in Appendix A).  $\slashed{\mathbf{1}}_n$ is the exchange matrix (the $(n\times n)$-matrix whose entries on the anti-diagonal are 1 with all other entries being 0). In the fourth column, the subscript $_{\text{left}}$ indicates the left regular representation of the subscripted algebra.  The not associative algebras in row 7 and row 8 are treated in Appendix B.}

\newpage

Table III collects various examples presented here and in Appendix A and Appendix B.   The special unital norm families, which can always be explicitly supplied for algebras of dimension less that five, are presented in the third column of Table III.
 This column makes it particularly easy to identify algebras of the same dimension that cannot be isomorphic, and algebras of different dimensions that cannot be related by an epimorphism (even if one is a subalgebra of another).  Corollary \ref{8/17/22.5} is also useful in this regard (e.g., consider the algebra in row 12 versus the algebra in row 13). This is accomplished without any computations related to identification of ideals.

Also in the third column of Table III, the factor in the expression for a special unital norm family that is devoid of an exponent employing a (Greek letter) parameter is the special unital norm family member that is closely allied with the usual norm of an algebra element.

\section{How large is the class of invariants implied by this approach?}

\subsection{Type I anti-rotors}
While we have derived a mechanism that generates \mbox{algebra} isomorphism invariants based on the specification of $f(s)=s^{-1}$ in Definition  
\ref{7/21/23.1}, there are infinitely many relevant possibilities for $f(s)$ dependent on the algebra's bilinear product that could substitute \mbox{for $s^{-1}$}.

\begin{definition}\label{10/29/25.1}  For $j\in\mathbb{Z}$, $j\ne 0,1$, if $s^j$ substitutes for $s^{-1}$ in Definition  
\ref{7/21/23.1}, the set of metrics satisfying the accordingly modified (\ref{5/22/25.3}) is the {\it Type I anti-rotor associated with} $f(s)=s^j$. \end{definition}

A trivial modification of the proof of Theorem \ref{6/6/23.1} results in,

\begin{corollary}\label{10/29/25.2} Given two isomorphic algebras, their respective Type I anti-rotors associated with $f(s)=s^{j}$ are the same up to a change of basis. \end{corollary} 

 One might then wonder if the association of an algebra with the above multiplicity of anti-rotors (including those based on linear combinations of the above $s^j$) could thereby be an endless source of algebra invariants.  If so, it would be appropriate to wonder how one could bring order to such a collection of processes that literally spew out invariants.    

But before contemplating this issue too deeply, we observe that
there is one thing more we can glean from the fourteen rows of algebras in Table III.  For each algebra, all of its Type I anti-rotors associated with the various $f(s)=s^j$ are identical for ``small" $|j|$.  Extrapolating, for any given algebra, if the above holds for {\it all} $j\in\mathbb{Z}$ except for $j= 0$ or $j=1$ (by definition), linearity of the exterior derivative means it would hold for arbitrary linear combinations of the various $s^j$ - implying that with respect to suitable domains in $\mathbb{R}^n$ it would hold for functions expressible by infinite series in these powers of $s$.  That is, the kernel of $\slashed{\text{S}}_f$ would be \mbox{independent of $f$}.

However, the behavior cited in the second sentence of the prior paragraph was unanticipated and seems surprising. There are several obvious avenues one might initially think would lead to an understanding of it.

First, let $\hat{s}$ be the left regular representation of an algebra with elements $s$, so that we can write $s^j = \hat{s}^{j-1}[s]$, where $s$ ranges over the algebra elements in an appropriate neighborhood of $\mathbf{1}$.  The behavior noted in the penultimate paragraph suggests that since $Ls$ is irrotational, and assuming $(L\hat{s}L^{-1})Ls = Ls^2$ is irrotational, then repeated applications of $L\hat{s}L^{-1}$ will continue to produce irrotational vector fields, i.e., $Ls^j = (L\hat{s}L^{-1})^{j-1}Ls = Ls^j$ is irrotational.  Why?
One might suspect that this follows because, in general, when an operator applied to an irrotational vector field results in an irrotational vector field, then it will continue to produce irrotational vector fields on all successive applications - which is false.  That is, consider the vector field $(F(s))[s]$ where $F(s)$ is an $(n\times n)$-matrix whose entries are smooth functions of components of $s$, and suppose $(F(s))[s]$ is irrotational.  This does {\it not} imply that $(F(s))^j[s]$ is irrotational for all $j\in\mathbb{Z}$. For example, let \mbox{$s=\left(\begin{array}{ccc}
 x   \\ 
y    
\end{array}\right)$,} and \mbox{$F(x,y)= \left(\begin{array}{ccc}
x  & 2y   \\ 
 2y & 2x    
\end{array}\right)$.}  Then it is easily checked that $(F(x,y))\left(\begin{array}{ccc}
 x   \\ 
y    
\end{array}\right)$ is \mbox{irrotational} and, in fact, $(F(x,y))^2\left(\begin{array}{ccc}
 x   \\ 
y    
\end{array}\right)$ is also irrotational.  But $(F(x,y))^3\left(\begin{array}{ccc}
 x   \\ 
y    
\end{array}\right)$ is not irrotational.  

Another possibility for understanding the behavior cited in the second paragraph following Corollary \ref{10/29/25.2}, is inspired by the easily appreciated truth of the generalization of that behavior to all $j\in\mathbb{Z}$ in the particular case of the full matrix algebra $M_n(\mathbb{R})$.  
Specifically, there is the well known Matrix Calculus expression valid for all $j\in\mathbb{Z}$ and for $s$ in a neighborhood of $\mathbf{1}$, \begin{equation}\label{9/28/25.1} \frac{\partial}{\partial s}\left(\text{Tr}(s^{j+1})\right) = (j+1)\left(s^j\right)^T,\end{equation} where $\text{Tr}$ indicates trace (in the Matrix Calculus literature, $s$ in (\ref{9/28/25.1}) is typically written as $\mathbf{X}$). Using Jacobi's formula, we have noted earlier that the transpose operation $^T$ (i.e., $\mathcal{T}$) is an uncurling metric for $M_n(\mathbb{R})$ in the exceptional $f(s)=s^{-1}$ case for which (\ref{9/28/25.1}) gives no useful information. Thus, for all other $j\in\mathbb{Z}$, (\ref{9/28/25.1}) indicates that $^T$ is indeed also a member (uncurling metric) of each of the Type I anti-rotors defined by the different $f(s)=s^j$ (since the expression on the left-hand-side of (\ref{9/28/25.1}) is a gradient).  So, according to an easy extrapolation of the argument in the proof of Proposition \ref{4/8/22.1} (extending the argument from $f(s)=s^{-1}$ to $f(s)=s^j$, $j\ne 0,1$, to again indicate that the only uncurling metrics for $M_n(\mathbb{R})$ are multiples of $\mathcal{T}$), here we have an example of an algebra whose space of uncurling metrics is indeed independent of the choice of $f(s)=s^j$, for $j\ne 0,1$.  On the other hand, it is not obvious how this observation can be adapted to the case of {\it sub}algebras of $M_n(\mathbb{R})$.  

As an aside, the observations in the last paragraph lead to,
\begin{proposition} \label{10/25/25.5}  For any given algebra, the Type I anti-rotor associated with $f(s)=s^j$ for any $j\in\mathbb{Z}$, $j\ne 0,1$, has nontrivial members.\end{proposition}
\begin{proof}
 According to the penultimate paragraph, the algebra $M_n(\mathbb{R})$ is such that its Type I anti-rotor associated with $f(s)=s^j$ for any $j\in\mathbb{Z}$ has nontrivial members - resulting from scalar multiples of the transpose operator. Identical arguments as appearing in the proofs of Lemma \ref{5/13/22.1} and Theorem \ref{11/17/22.1}, except replacing $s^{-1}$ with $s^j$, then establish the proposition.\end{proof}

Returning to our original quest, a third way to try to find an explanation for the behavior cited in the second paragraph following Corollary \ref{10/29/25.2}, is to exploit the fact that pullbacks commute with the exterior derivative.  Thus, suppose it could be shown that for any $j\ne 0,1$, $j\in\mathbb{Z}$, one can find a $j$-dependent diffeomorphism $\phi$  between relevant domains in $\mathbb{R}^n$ resulting in a $j$-dependent pullback $\phi^*$ between the 1-forms $(\mu s^{-1})\cdot \mathbf{d}s$ and $(\mu s^j)\cdot \mathbf{d}s$.   Then we would have $\phi^* \big[(\mu s^{-1})\cdot \mathbf{d}s\big] = (\mu s^j)\cdot \mathbf{d}s$, along with a similar equation expressing $\big(\mu s^{-1}\big)\cdot \mathbf{d}s$ as a pullback of $\big(\mu s^j\big)\cdot \mathbf{d}s$.  Commutativity of the exterior derivative and a pullback would then imply that the exterior derivative of $(\mu s^j)\cdot \mathbf{d}s$ vanishes if and only if the exterior derivative of $(\mu s^{-1})\cdot \mathbf{d}s$ vanishes - and we would have our explanation.

However, proving that all those pullbacks exist would involve a demonstration of the existence of solutions to each set of nonlinear partial differential equations arising from the choice of any given algebra.  For individual cases of a particular algebra and a particular value of $j$, along with the stipulation that the relevant diffeomorphisms must have only rational function components, the problem might be addressed with known techniques of computational algebraic geometry.  But it is difficult to overcome the impression that the general problem is intractable, i.e., that of establishing the existence of the requisite pullbacks for any algebra and any $j\ne 0,1$. 

Thus, at this point, our limited computational evidence only leaves us with,

\begin{conjecture}\label{10/30/25.1} All of the Type I anti-rotors of a given algebra are identical. \end{conjecture}

We have already noted that linearity of the exterior derivative means that the conjectured equivalence of anti-rotors pertains not just to the above powers of $s$, but also to arbitrary linear combinations of such powers of $s$ - implying that with respect to suitable domains in $\mathbb{R}^n$ it would hold for functions expressible by infinite series in these powers.  Thus, it would hold for the ``algebra-analytic" transformations from $\mathbb{R}^n$ to itself, i.e., vector fields having the form, \[ f(s) = \sum_{j=-\infty}^{\infty}\alpha_j s^j,\] so long as there is at least one term where $\alpha_j\ne 0$ for $j\ne 0,1$ (i.e., so long as $f(s)$ is dependent on the algebra's bilinear product).  

Hence, the above conjecture can be put in the following equivalent form.

\medskip

\noindent {\bf Conjecture 8.4$'$}.  Suppose $f(z)$ and $g(z)$ are two functions of a complex variable that are holomorphic in some annulus around the origin, all coefficients of their respective Laurent series are real, and each of these series have at least one term $\alpha_j z^j$ where $\alpha_j\ne 0$ for $j\ne 0,1$.  Let $q$-dimensional algebra $A$ be a subalgebra of $M_n(\mathbb{R})$.  Consider a domain $D$ in $A$ where all members of \mbox{$\{s: s\in D\}\cup \{s^{-1}: s\in D\}$} are such that their eigenvalues lie in the above annulus.   Then with respect to $D$, and any given $\mu\in \text{Sym}_q$, $\mu f(s)$ is irrotational if and only if $\mu g(s)$ is irrotational, where $f(s)$ and $g(s)$ are treated as $q$-dimensional vector fields. 

\medskip

In essence, recalling Section 2, the conjecture asserts that the kernel of $\slashed{\text{S}}_f$ is independent of $f$ over the relevant class of functions that depend on the algebra product. 

To summarize: On its face, it would seem very unlikely that Conjecture \ref{10/30/25.1} is true.  And yet, one can generate a great deal of computational evidence in its favor - beginning with the algebras listed in Table III.   But either way, resolution of the conjecture will lead to an interesting result.
\begin{itemize}\item If Conjecture \ref{10/30/25.1} is true, then every algebra is uniquely associated with an anti-rotor (all its Type I anti-rotors are actually the same).  This would be a nice outcome for what initially seemed to be a disorderly situation.  Moreover, the truth of Conjecture \ref{10/30/25.1} might be a deep result, as it seems so counter-intuitive.  \item
If Conjecture \ref{10/30/25.1} is false, the set of invariants implied by our program will be (perhaps much) greater than we presently suspect, and new organizing principles will likely need to be developed to bring order and understanding to the full set of invariants generated by this program.
\end{itemize} 

However, the above is not the end of our current story.  As we shall see in the next subsection, for any particular choice of algebra, we can at least identify infinitely many different infinite sequences of vector fields $s^j$, such that within a given sequence each of its members is uncurled by the same ``Type II anti-rotor", where each Type II anti-rotor is a subset of a Type I anti-rotor.

\subsection{Type II anti-rotors}
We have defined the Type I anti-rotors based on different choices of powers of $s$ replacing $s^{-1}$ in Definition \ref{7/21/23.1}.  We now introduce ``Type II anti-rotors", which will be generally referred to individually as ``$k$th-order anti-rotors".  Each of these is a set of metrics that uncurls infinitely many vector fields given by certain powers of $s$.  They are defined in a completely different manner than Type I anti-rotors.
   
Consider the 1-form $(\mu s)\cdot \mathbf{d}s$ for $\mu\in \text{Sym}_n$.    The mapping $\phi(s) = s^2$ is a diffeomorphism between appropriate domains in $\mathbb{R}^n$.  This generates a pullback which operates on the 1-form in the penultimate sentence according to, \begin{equation}\label{10/24/25.1} \phi^* \left[(\mu s)\cdot\mathbf{d}s\right] =  (\mu s^2)\cdot \mathbf{d}(s^2) = (\mu s^2)\cdot (s\,\mathbf{d}s + (\mathbf{d}s)s) = (\mu s^2)\cdot ((\hat{s}_L+\hat{s}_R)[\mathbf{d}s]),\end{equation} where $\hat{s}_L$ is the left regular representation, and $\hat{s}_R$ is the matrix effecting right multiplication in the given algebra.  Suppose we can find nontrivial $\mu'\in \text{Sym}_n$ (i.e., $\mu'$ not the matrix with all entries being 0) such that, \begin{equation}\label{10/24/25.3}(\hat{s}^T_L+\hat{s}_R^T)\mu' = \mu'(\hat{s}_L+\hat{s}_R).\end{equation}   Satisfaction of (\ref{10/24/25.3}) implies, \begin{equation}\label{10/24/25.2} \phi^* \left[(\mu' s)\cdot\mathbf{d}s\right]  =\left((\hat{s}^T_L+\hat{s}_R^T) [\mu' s^2]\right)\cdot\mathbf{d}s =  (\mu'(\hat{s}_L+\hat{s}_R)[s^2])\cdot\mathbf{d}s = 2\, (\mu' s^3) \cdot \mathbf{d}s,\end{equation}  where in obtaining the first equality above from (\ref{10/24/25.1}) we have used the characteristic property of an adjoint in the context of an inner product.  Note that the exterior derivative of the 1-form operated on by $\phi^*$ on the left-hand-side of (\ref{10/24/25.2}) is zero.  Since pullbacks commute with the exterior deriviative, (\ref{10/24/25.2}) implies that $\mu's^3$ is irrotational, since $\mu' s$ is irrotational (note that for $\mu\in \text{Sym}_n$, $\mu s$ is always irrotational, but in general $\mu s^3$ will not be irrotational, since (\ref{10/24/25.3}) is required to effect (\ref{10/24/25.2})).  In fact, we can proceed indefinitely in this manner (i.e., next replacing $\mu' s$ on the left-hand-side of (\ref{10/24/25.2}) with $\mu' s^3$, etc.) to show that (\ref{10/24/25.3}) implies that $\mu's^j$ is irrotational, where $j$ is any member of the sequence defined by $a_1=3$ and $a_{q+1} = 2a_q +1$, i.e., the sequence $(3, 7, 15, 31, \dots, 2^{q+1} -1,\dots)$.  So, all of the sudden, we have found an infinite sequence of powers of $s$ (with exponentially increasing spacing between consecutive powers) that are all uncurled by the same metric $\mu'$ based on the latter's satisfaction of (\ref{10/24/25.3}).  In fact, again representing $\text{Sym}_n$ by the symmetric parametrized matrix $M$ as defined in Section 2, we can compute the {\it collection} of metrics satisfying (\ref{10/24/25.3}) and term  it the ``Type II second-order anti-rotor", which is a parametrized matrix that uncurls the all vector fields in the sequence $\{f_{a_q}(s)=s^{a_q}:q\in\mathbb{N}\}$ (``second order", because the exponent in the utilized diffeomorphism $\phi(s)=s^2$ is 2).  We could specify that the first-order anti-rotor is associated with the trivial diffeomorphism $\phi(s)=s$, thereby being the set of metrics for which $\mu' s$ has zero curl.  Thus, the first-order Type II anti-rotor would be simply $\text{Sym}_n$.  However, we will not recognize it as a Type II anti-rotor.  That is, our Type II anti-rotors begin with the second order.

Proceeding analogously beyond the second-order anti-rotor, we can define the $k$th-order anti-rotor.  Thus, instead of confining ourselves to the diffeomorphism $\phi(s)=s^2$ as utilized in (\ref{10/24/25.1}) and (\ref{10/24/25.3}), we utilize the diffeomorphism $\phi(s)=s^k$, for any given choice of $k\in\mathbb{N}$, $k>1$.   Accordingly, we obtain
\begin{multline}\label{10/27/25.1} \phi^*[(\mu' s)\cdot \mathbf{d}s] = (\mu' s^k)\cdot \mathbf{d}s^k =  (\mu' s^k)\cdot \left(\sum_{q=0}^{k-1} s^q\, \mathbf{d}s\, s^{k-q-1}\right) \\ = (\mu' s^k)\cdot \left(\left(\sum_{q=0}^{k-1} \hat{s}_L^q\hat{s}_R^{k-q-1}\right) [\mathbf{d}s]\right) =\left(\left(\sum_{q=0}^{k-1} \hat{s}_L^q\hat{s}_R^{k-q-1}\right)^T \mu' [s^k]\right)\cdot  \mathbf{d}s \\ = 
\left((\mu' s^k)\left(\sum_{q=0}^{k-1} \hat{s}_L^q\hat{s}_R^{k-q-1}\right)[s^k]\right)\cdot \mathbf{d}s = k(\mu' s^{2k-1})\cdot \mathbf{d}s,\end{multline} where the fifth equality requires that $\mu'$ satisfies a more complicated equation analogous to (\ref{10/24/25.3}), i.e.,  in place of the term $(s\,\mathbf{d}s + (\mathbf{d}s)s)$ that appears in (\ref{10/24/25.1}).  Specifically, the $k$th-order anti-rotor is the set of $\mu'$ satisfying, \begin{equation}\label{10/29/25.5} \left(\sum_{q=0}^{k-1} \hat{s}_L^q\hat{s}_R^{k-q-1}\right)^T \mu' = 
\mu' \sum_{q=0}^{k-1} \hat{s}_L^q\hat{s}_R^{k-q-1}.\end{equation}
As with the $k=2$ case in the prior paragraph, we can insert $s^{2k-1}$ on the right-hand-side of (\ref{10/27/25.1}), etc., to obtain the sequence of vector fields $s^j$ that are each uncurled by any of the members of this $k$th-order anti-rotor.  In this way we obtain infinitely many different infinite sequences of vector fields, every member of the sequence being uncurled by the same $k$th-order anti-rotor.  Analogous to the second-order case above, the members of the $k$th-order anti-rotor will uncurl those vector fields $s^j$ with $j$ a member of the sequence defined by $a_1=2k-1$, $a_{q+1} = k\, a_q + k-1$.   
 
 Based on (\ref{10/29/25.5}), we formalize this as,
\begin{definition}\label{10/29/25.3} Let $\mathcal{D}_k$ be the operator which when applied to $\mathbf{d}s$ results in $\mathbf{d}s^k$, i.e., given by the sum on the right-hand-side of (\ref{10/29/25.5}).   For $k>1$, the {\it $k$th-order Type II anti-rotor} is the space of metrics each of whose members results in a symmetric operator when it left-multiplies $\mathcal{D}_k$.\end{definition}

\begin{proposition}\label{10/29/25.4} A $k$th-order Type II anti-rotor is a subset of each of the Type I anti-rotors associated with $f(s)=s^j$, where $j$ is any member of the sequence defined by $a_1=2k-1$, $a_{q+1} = k\, a_q +k-1$. \end{proposition}
\begin{proof} The Type I anti-rotor associated with a particular $f(s)=s^j$ encompasses all metrics uncurling the latter vector field.\end{proof} 

In view of Corollary \ref{10/29/25.2}, the above proposition also indicates that Type II anti-rotors associated with isomorphic algebras are the same up to a change of basis.

To put things in perspective, on the one hand Proposition \ref{10/25/25.5} indicates that for a given algebra, its Type I anti-rotor associated with any $f(s)=s^j$, $j\ne 0,1$, has nontrivial members - although the relationships between all these Type I anti-rotors are as yet undetermined (indeed, based on the computational evidence cited above, our problem is to decide whether or not they are all identical).  On the other hand, we have the $k$th-order Type II anti-rotors defined by pullbacks, which for each $k$ lead to an infinite collection of vector fields whose members are powers of $s$ that are uncurled by the {\it same} parametrized matrix of uncurling metrics.  However, we have not shown that Type II anti-rotors necessarily have any nontrivial members, we have no general understanding of the connections between the Type II anti-rotors of different orders, and so far we know little more than that they are subsets of Type I anti-rotors.

But there is more to consider.

For any given odd number $j>1$, Proposition \ref{10/29/25.4} indicates that there is a $k$th-order Type II anti-rotor that addresses the vector field $s^j$.  But the proposition says nothing about the case where $j$ is an even number.  Therefore, in the context of the diffeomorphism $\phi(s)=s^2$ between appropriate domains, let us examine $\phi^* \left[(\mu s^2)\cdot\mathbf{d}s\right] =  (\mu s^4)\cdot \mathbf{d}(s^2)$.  If $\mu'$ is a member of the {\it second-order} Type II anti-rotor, we have  
\begin{eqnarray} \nonumber \phi^* \left[(\mu' s^2)\cdot\mathbf{d}s\right]  &=& (\mu' s^4)\cdot (s\mathbf{d}s+\mathbf{d}s) = \left((\hat{s}^T_L+\hat{s}_R^T) [\mu' s^4]\right)\cdot\mathbf{d}s \\  &=& (\mu'(\hat{s}_L+\hat{s}_R)[s^4])\cdot\mathbf{d}s =  2\, (\mu' s^5) \cdot \mathbf{d}s. \label{10/30/25.5}\end{eqnarray} 
From Proposition \ref{10/29/25.4}, we know that $s^5$ is uncurled by the {\it third-order} anti-rotor members.  But now let us assume that all orders of Type II anti-rotors are identical.  In that case, $\mu'$ also uncurls $s^5$.  From the commutativity of the exterior derivative and a pullback, it would then follow from (\ref{10/30/25.5}) that $\mu'$ uncurls $s^2$.  We can repeat this type of argument for all $s^j$ with $j$ a positive even number.  Thus, equivalence of Type II anti-rotors of all orders would imply that the unique Type II anti-rotor uncurls all vector fields $s^j$, $j\in\mathbb{N}$, $j>1$.

To show that all orders of Type II anti-rotors are identical, according to Definition \ref{10/29/25.3}, it is ``only" necessary to show that all $\mathcal{D}_k$ are symmetrized by being left-multiplied by the same parametrized matrix.  Unfortunately, we don't have a proof of that assertion either, so the computation evidence cited below will only leave us with,

\begin{conjecture}\label{10/30/25.6} All of the Type II anti-rotors of a given algebra are identical. \end{conjecture}

 Clearly, our program would be greatly simplified if there is a unique Type I anti-rotor, and there is a unique Type II anti-rotor, and these are equal to each other.  Although those possibilities might seem fanciful, consider the following sample of available computational evidence.  


As with the Type I anti-rotors, we can use the algebras in Table III to examine Type II anti-rotors.  In particular, it is easy to show that for every associative algebra in Table III, its second-order Type II anti-rotor is identical to its Type I anti-rotor associated with $f(s)=s^{-1}$ - as well being identical to the Type I anti-rotors associated with \mbox{$f(s)=s^j$}, for small $|j|$, $j\in \mathbb{Z}$, $j\ne 0,1$, e.g., with $f(s)=s^3$.   
For example, we can illustrate this with the noncommutative algebra in the twelfth row of \mbox{Table III} (being noncommutative, we have, $\hat{s}_L\ne \hat{s}_R$).  For this algebra, we obtain,
\[ \hat{s}_L = \left[
\begin{array}{cccc}
x & 0 & 0 & 0   \\
y & x & 0 & 0 \\
z & 0 &  x & 0   \\
w & 0 &  y & x
\end{array}
\right], \hspace{5mm} \hat{s}_R = \left[
\begin{array}{cccc}
x & 0 & 0 & 0   \\
y & x & 0 & 0 \\
z & 0 &  x & 0   \\
w & z &  0 & x  
\end{array}\right], \hspace{5mm} M_\mathfrak{u} = \left[
\begin{array}{cccc}
\alpha & \beta & \gamma & \delta   \\
\beta & 0 & \frac{1}{2}\delta & 0 \\
\gamma & \frac{1}{2}\delta &  0 & 0   \\
\delta & 0 &  0 & 0
\end{array}
\right] ,\] where $M_\mathfrak{u}$ is the Type I anti-rotor associated with $f(s)=s^{-1}$ (i.e., the anti-rotor as forwarded by Definition \ref{7/21/23.1}), which is the same as the Type I anti-rotor associated with $f(s)=s^3$.
Equation (\ref{10/24/25.3}) is then easily verified using the parametrized $M_\mathfrak{u}$ as a substitute for $\mu'$ - which means that the second-order Type II anti-rotor is identical to $M_\mathfrak{u}$, since the second-order Type II anti-rotor is a subset of a corresponding Type I anti-rotor, and in this example we have shown that every member of the Type I anti-rotor associated with $f(s)=s^{-1}$ (or $f(s)=s^3$) is also a member of the second-order anti-rotor (i.e., it satisfies (\ref{10/24/25.3})).  Without much effort, using the more general (\ref{10/29/25.5}) with $k=3$, it is seen that this also holds for the third-order anti-rotor, i.e., the third-order anti-rotor is equal to $M_\mathfrak{u}$.   

Indeed, we can easily generate computational evidence suggesting, 

\begin{conjecture}\label{10/28/25.1} For a given algebra, all Type I anti-rotors are identical, all Type II anti-rotors are identical, and the unique Type I anti-rotor is identical to the unique Type II anti-rotor. \end{conjecture}

If the above conjecture is true, every algebra is uniquely associated with an anti-rotor (in the language of Section 1, the kernel of the Skewer is independent of $f$), and we have the class of invariants introduced in the Sections 1 through 7.  On the other hand, if any of the three conjectures in this section is false, we will still have the aforementioned class of invariants, but there will apparently be much additional interesting algebra-dependent structure to be explored.
  
  

\vspace{.75cm}
 
 
 
\appendix

\section{Upper triangular Toeplitz matrix algebras} 

The anti-rotor and unital norm group associated with an upper triangular Toeplitz matrix algebra have some interesting features.
\begin{itemize}\item The anti-rotor is given by the left regular representation (the algebra whose elements are lower triangular Toeplitz matrices) multiplied on the left by the exchange matrix. \item The usual norm is sensitive to only a single component of an element, but almost all unital norms are sensitive to all components of an element.\end{itemize}


 \begin{definition}\label{5/4/22.1} $\mathbb{T}_n$ is the $n$-dimensional vector space of real upper triangular Toeplitz $(n\times n)$-matrices. For $s\in \mathbb{T}_n$, the notation $s=(x_1,\dots,x_n)$ indicates that $s$ corresponds to the Toeplitz upper triangular $(n\times n)$-matrix whose entries on the successive diagonals, beginning with the main diagonal and proceeding to the right, are given by the successive vector component values $x_1,\dots,x_n$.   
\end{definition}

In the sequel, for $s\in \mathbb{T}_n$, $s$ can refer either to its format as an $n$-vector or its format as an $(n\times n)$-matrix, depending on the context.


The following result is known, but we include the proof for completeness, and to introduce some notation.  



 \begin{proposition}\label{5/10/22.1} Under the usual matrix product, $\mathbb{T}_n$ is a real unital commutative associative $n$-dimensional algebra.\end{proposition} 
 
 \begin{proof}
Let ${\bf z}$ be the $(n \times n)$-matrix with entries having value $1$ on its superdiagonal  (the diagonal immediately to the right of the main diagonal) and value $0$ otherwise.  It is easily shown that for $j<n$ its $j$-th power, ${\bf z}^j$, is the $(n\times n)$-matrix with entries $1$ on the $j$-th diagonal to the right of the main diagonal, with all other entries being $0$.   Thus, we can express the matrix format of $s\in\mathbb{T}_n$ as a polynomial in ${\bf z}$, such that,
\begin{equation}\label{5/11/22.5} s = x_1 {\bf 1} + x_2 {\bf z} + x_3 {\bf z}^2 + \cdots + x_n {\bf z}^{n-1}.\end{equation}  It is similarly easy to show that ${\bf z}^k$ is the matrix of all zero entries if $k\ge n$.  It follows that for $s,t\in \mathbb{T}_n$, we have $st\in \mathbb{T}_n$.   Based on the polynomial format of an element as given in (\ref{5/11/22.5}), the proposition follows.
 \end{proof}
 

We will show that upper triangular Toeplitz matrix algebras have unital norms sensitive to all element components, providing a great contrast with the usual norm for these algebras, which is sensitive to only one element component.  One can easily generate other algebras with unital norms ``similarly dominating" the usual norm, as seen in some other examples in Section 7.

To approach the problem of deriving the anti-rotor for $\mathbb{T}_n$, the following will be important.

\begin{lemma} \label{5/7/22.15} For a unit $s\in \mathbb{T}_n$,  $n>1$, let $\left( s^{-1} \right)_p$ be the $p$-th component of $s^{-1}$ in the latter's $n$-vector format.  Given $k\le n$, then for all $j<k$ we have \begin{equation}\label{5/6/22.10} \frac{\partial \left( s^{-1} \right)_k}{\partial x_j} = \frac{\partial \left( s^{-1} \right)_{k-j+1}}{\partial x_1}.\end{equation}
\end{lemma}

\begin{proof} The proof is by induction. Equation (\ref{5/6/22.10}) is a tautology for the $k=2$ case.  We thus assume the inductive hypothesis that (\ref{5/6/22.10}) is true  for all cases up through $k-1$, and then establish its truth for the $k$-th case.


We must first express $s^{-1}$ in terms of the components of  $s=(x_1,\dots,x_n)\in \mathbb{T}_n$ with $x_1\ne 0$.   From Proposition \ref{5/10/22.1},  $s^{-1}$ is an upper triangular Toeplitz matrix.  But as an $n$-vector, it is easily seen from (\ref{5/11/22.5}) that we can write it as  \begin{equation}\label{5/9/22.1} s^{-1} = (\frac{1}{x_1}, \frac{y_1}{x_1},\dots, \frac{y_{n-1}}{x_1}),\end{equation} where the $y_i$ are to be determined. 
Based on the notation introduced in the proof of Proposition \ref{5/10/22.1}, we can use (\ref{5/11/22.5}) to write 
\begin{align}\nonumber s &= x_1 \left({\bf 1} + \frac{x_2}{x_1} {\bf z} + \frac{x_3}{x_1} {\bf z}^2 + \cdots + \frac{x_n}{x_1} {\bf z}^{n-1}\right), \\ s^{-1} &= x_1^{-1} \left({\bf 1} + y_1 {\bf z} + y_2 {\bf z}^2 + \cdots + y_{n-1} {\bf z}^{n-1}\right).  \label{5/18/22.10} \end{align}  Substituting the above two equations into $s^{-1} s = {\bf 1}$ leads to following set of relations, which allow the $y_i$ to be computed recursively:
\begin{align} -y_1 &= \frac{x_2}{x_1} \nonumber \\ -y_2 &= \frac{x_2}{x_1}y_1 + \frac{x_3}{x_1} \nonumber \\ &\;\vdots \nonumber \\ -y_{k-1} &= \frac{x_2}{x_1} y_{k-2} + \frac{x_3}{x_1} y_{k-3} + \cdots + \frac{x_k}{x_1} \label{5/7/22.1} \\ &\;\vdots \nonumber \end{align} 
 
 \noindent For $k>1$, (\ref{5/18/22.10}) implies  $\left( s^{-1} \right)_k = \frac{y_{k-1}}{x_1}$.  Applying this to (\ref{5/7/22.1}) with $k>1$, we find
\begin{eqnarray} \nonumber \left( s^{-1} \right)_k &=& -\frac{x_2}{x_1}\left( s^{-1} \right)_{k-1} - \frac{x_3}{x_1}\left( s^{-1} \right)_{k-2} -\dots  - \frac{x_k}{x_1}\left( s^{-1} \right)_1 \\ &=& -\frac{1}{x_1}\sum_{m=1}^{k-1} x_{m+1}\left( s^{-1} \right)_{k-m}. \label{5/7/22.2}\end{eqnarray}
Differentiating with respect to $x_j$, we find \begin{equation}  \frac{\partial \left( s^{-1} \right)_k}{\partial x_j} = -\frac{1}{x_1} \sum_{m=1}^{k-j} x_{m+1} \frac{\partial \left( s^{-1} \right)_{k-m}}{\partial x_j} -  \frac{1}{x_1}\left( s^{-1} \right)_{k-j+1} \label{5/7/22.4}, \end{equation} where the upper limit on the sum on the right-hand-side is $(k-j)$ rather than $(k-1)$ because the partial derivatives in the sum are zero for $m> k-j$ (i.e., $\left( s^{-1} \right)_p$ is not a function of $x_j$ for $p<j$).  

Next, substituting $(k-j+1)$ for $k$ in (\ref{5/7/22.2}) and then differentiating with respect to $x_1$, we have   \begin{eqnarray} \nonumber
 \frac{\partial \left( s^{-1} \right)_{k-j+1}}{\partial x_1} &=& \frac{1}{x_1^2} \sum_{m=1}^{k-j} x_{m+1} \left( s^{-1} \right)_{k-j-m+1}    - \frac{1}{x_1} \sum_{m=1}^{k-j} x_{m+1} \frac{\partial \left( s^{-1} \right)_{k-j-m+1}}{\partial x_1}  \\ &=& - \frac{1}{x_1}\left( s^{-1} \right)_{k-j+1}  -  \frac{1}{x_1} \sum_{m=1}^{k-j} x_{m+1} \frac{\partial \left( s^{-1} \right)_{k-j-m+1}}{\partial x_1},\label{5/7/22.3}\end{eqnarray} where the equality of the first terms on the right-hand-sides of the two above equalities comes from (\ref{5/7/22.2}) itself (with $(k-j+1)$ substituting for $k$). 
 Now, according to the induction hypothesis, $\frac{\partial \left( s^{-1} \right)_{k-m}}{\partial x_j} = \frac{\partial \left( s^{-1} \right)_{k-j-m+1}}{\partial x_1}$.  Inserting this into (\ref{5/7/22.3}), we see that the left-hand-sides of (\ref{5/7/22.4}) and (\ref{5/7/22.3}) are equal, as was required to be shown.
 \end{proof} 

Let $\slashed{\bf 1}_n$ be the $(n\times n)$ exchange matrix (the matrix whose entries are 1 on the antidiagonal and zero everywhere else).

\begin{theorem} \label{5/4/22.2} For any $t\in \mathbb{T}_n$ let $\hat{t}$ be its $(n\times n)$-matrix format.    Then $\hat{t} \slashed{\bf 1}_n$ is an uncurling metric.  Conversely, each member of the anti-rotor associated with $\mathbb{T}_n$ can be expressed as $\hat{t}\slashed{\bf 1}_n$ for some $t\in \mathbb{T}_n$.   \end{theorem}


\begin{proof} The following notation is more convenient for the moment.  Given $(\gamma_n,\dots,\gamma_1)$ $\in \mathbb{T}_n$, let $G_n$ be its upper triangular Toeplitz $(n\times n)$-matrix format (which happens to be the transpose of its left regular representation), and consider the resulting Hankel matrix, 
\begin{equation}\label{5/7/22.20} H_n \equiv G_n \slashed{\bf 1}_n = \begin{bmatrix} \gamma_1 & \gamma_2 & \gamma_3 & \dots & \gamma_n \\ \gamma_2 & \gamma_3 & \dots & \gamma_n & 0 \\ \gamma_3 & \dots & \iddots & \iddots & 0  \\ \vdots & \iddots & \iddots & \iddots & \vdots \\ \gamma_n & 0 & \dots & \ldots & 0 \end{bmatrix}.\end{equation}
For $\left( s^{-1} \right)_p$ as defined in the statement of Lemma \ref{5/7/22.15}, we then have
\begin{multline}  H_n s^{-1}   =  \Big( \gamma_1\left( s^{-1} \right)_1  + \dots + \gamma_n \left( s^{-1} \right)_n, \\ \gamma_2\left( s^{-1} \right)_1  + \dots + \gamma_n \left( s^{-1} \right)_{n-1}, 
  \dots , 
 \gamma_n \left( s^{-1} \right)_1 \Big). \label{5/7/22.40}
\end{multline}
 We will first demonstrate that $H_n$ is an uncurling metric.  According to (\ref{5/22/25.3}) as referred to by Definition \ref{5/22/25.3}, we are required to demonstrate 
\begin{equation}\label{5/7/22.30} d\left([H_n s^{-1}]\cdot \mathbf{d}s\right) = 0.\end{equation}  

To verify (\ref{5/7/22.30}), we apply an induction argument.  The $n=2$ case, pertaining to the dual numbers, is easily demonstrated to satisfy the above equation.  We assume the induction hypothesis that (\ref{5/7/22.30}) holds for all cases less than $n$, and now need to show that it holds for the $n$-th case.  
We begin by showing that the partial derivative of the first component of $H_n s^{-1}$ with respect to any $x_j$, with $j>1$, minus the partial derivative of the $j$-th component of $H_n s^{-1}$ with respect to $x_1$, is zero. We will then show that the analogous difference of partial derivatives of {\it any} two different components of $H_ns^{-1}$ is also zero.

So, from (\ref{5/7/22.40}), 
\begin{equation}\label{5/18/22.5} \frac{\partial \left(H_n s^{-1}\right)_j}{\partial x_1} = 
\sum_{i=j}^n \gamma_i\frac{\partial \left(s^{-1}\right)_{i-j+1}}{\partial x_1}. 
\end{equation} 
But, \begin{equation}\label{5/18/22.6}  \frac{\partial \left(H_n s^{-1}\right)_1}{\partial x_j} = \sum_{i=1}^n \gamma_i\frac{\partial \left(s^{-1}\right)_i}{\partial x_j}
= \sum_{i=j}^n \gamma_i\frac{\partial \left(s^{-1}\right)_i}{\partial x_j} = \sum_{i=j}^n \gamma_i\frac{\partial \left(s^{-1}\right)_{i-j+1}}{\partial x_1},
\end{equation} where the second equality follows (\ref{5/9/22.1}) and (\ref{5/7/22.1}), which indicate that $\frac{\partial \left(s^{-1}\right)_i}{\partial x_j} = 0$ for $i<j$, and the third equality follows from Lemma \ref{5/7/22.15}.
Equations (\ref{5/18/22.5}) and (\ref{5/18/22.6})  demonstrate that the partial derivative of the first component of $H_n s^{-1}$ with respect to any $x_j$, with $j>1$, minus the partial derivative of the $j$-th component of $H_n s^{-1}$ with respect to $x_1$, is equal to zero.

Now we will consider the vanishing of the partial derivative of the $k$-th component of $H_n s^{-1}$ with respect to $x_j$ minus the partial derivative of the $j$-th component of $H_n s^{-1}$ with respect to $x_k$, for all possible $k>j>1$ (i.e, all choices relevant to (\ref{5/7/22.30}) other than considered in the prior paragraph).  In fact, based on the telescoping format of (\ref{5/7/22.40}), we already have the result that this difference vanishes, because it merely reduces to an instance of the theorem for the case of a matrix of the form $H_ks^{-1}$ with $k<n$ - and this case is assumed to hold according to the induction hypothesis.  So, (\ref{5/7/22.30}) is established.

Furthermore, a Hankel matrix of the form (\ref{5/7/22.20}) is the only possibility for an uncurling metric $L$ because any other choice resulting in the coefficients of the terms in any two components of $Ls^{-1}$ cannot lead to a null value for the difference of the two requisite partial derivatives (as is necessary for $d\left([L s^{-1}]\cdot \mathbf{d}s\right) = 0$) due to linear independence of the individual terms summing to the components. That is, replacing  $H_n$ with matrix $L=(a_{p,q})_{p,q=1}^n$ in (\ref{5/18/22.5}) 
and (\ref{5/18/22.6}), and invoking Lemma \ref{5/7/22.15}, we obtain
\begin{eqnarray}\nonumber \frac{\partial \left(L s^{-1}\right)_j}{\partial x_1}  -  \frac{\partial \left(L s^{-1}\right)_1}{\partial x_j} &=&  \sum_{i=1}^n a_{j,i}\frac{\partial \left(s^{-1}\right)_i}{\partial x_1} - \sum_{i=j}^n a_{1,i}\frac{\partial \left(s^{-1}\right)_i}{\partial x_j} \\  &=& \sum_{i=1}^n a_{j,i}\frac{\partial \left(s^{-1}\right)_i}{\partial x_1} - \sum_{i=j}^n a_{1,i}\frac{\partial \left(s^{-1}\right)_{i-j+1}}{\partial x_1} \nonumber \\
&=& \sum_{i=1}^n a_{j,i}\frac{\partial \left(s^{-1}\right)_i}{\partial x_1} - \sum_{i=1}^{n-j+1} a_{1,i+j-1}\frac{\partial \left(s^{-1}\right)_i}{\partial x_1}.
\label{5/9/22.10}
\end{eqnarray}
This indicates that we must have $a_{j,i}=0$ for $i>n-j+1$, and otherwise have $a_{j,i} = a_{1,i+j-1}$.  Since this must hold for any $j>1$, $L$ must be a Hankel matrix of the form (\ref{5/7/22.20}).  Thus, the $H_n$ are the only possibilities for uncurling metrics. 

Since $t=(\gamma_n,\dots,\gamma_1)$ is an arbitrary member of $\mathbb{T}_n$, the theorem follows.
\end{proof}

\begin{corollary}\label{7/1/22.1} For any $t\in \mathbb{T}_n$, denote the value of the $(1,n)$ entry of its matrix format $\hat{t}$ as $\hat{t}_{1,n}$.  Then $\frac{1}{\hat{t}_{1,n}}\hat{t}\slashed{\bf 1}_n$ is a normalized uncurling metric if $\hat{t}_{1,n}\ne 0$, and all normalized uncurling metrics can be written in this form for some $t\in \mathbb{T}_n$. \end{corollary}

\begin{proof} Theorem \ref{5/4/22.2} has shown that all uncurling metrics of $\mathbb{T}_n$ have the form (\ref{5/7/22.20}).
We now consider constraint (\ref{7/7/21.1}) of Definition \ref{5/22/25.3}.  Let $H_n^{(1)}$ be matrix $H_n$ except with the  $(1,1)$-entry having the value $1$.  
Since $\|{\bf 1}\|^2 =1$ for this algebra, we must show that \[ s^T H_n^{(1)} s^{-1} = s\cdot [H_n^{(1)}s^{-1}] = 1.\] From (\ref{5/9/22.1}) and (\ref{5/7/22.20}), we have
\begin{eqnarray}\nonumber s^T H_n s^{-1} &=& (x_1,\dots,x_n)\cdot \left[ H_n \left(\frac{1}{x_1}, \frac{y_1}{x_1},\dots,\frac{y_{n-1}}{x_1}\right)^T\right] \\ 
 &=&  \gamma_1 + \gamma_2y_1 + \gamma_3y_2 \hspace{3mm} + \hspace{3mm}\gamma_4y_3 \hspace{3mm} + \dots   \hspace{6mm}+\gamma_ny_{n-1} \nonumber \\
&+&  \gamma_2 \frac{x_2}{x_1} + \gamma_3\frac{x_2}{x_1}y_1 + \hspace{2mm} \gamma_4\frac{x_2}{x_1}y_2 + \dots +  \gamma_n\frac{x_2}{x_1}y_{n-2} \nonumber \\ &+&\gamma_3 \frac{x_3}{x_1}  + \hspace{3mm}\gamma_4\frac{x_3}{x_1}y_1 +  \dots \hspace{4mm} +  \gamma_n\frac{x_3}{x_1}y_{n-3} \nonumber \\ 
&+&  \hspace{5mm} \vdots \hspace{2.5cm} \iddots \bigskip\bigskip \nonumber \\ &+& \gamma_{n-1} \frac{x_2}{x_1} \hspace{3mm} + \gamma_n \frac{x_2}{x_1}y_1 \nonumber \\ &+& \hspace{2mm} \gamma_n\frac{x_n}{x_1}. \label{6/23/22.1}
 \end{eqnarray}
 For any $i>1$, if we collect all the terms on the right-hand-side above that each have $\gamma_i$ as a coefficient, and then apply (\ref{5/7/22.1}), it is quickly seen that the sum of all terms having $\gamma_i$ as a coefficient vanishes.  It follows that
 $s^T H_n s^{-1} =\gamma_1$.   Thus, the requisite matrix satisfying constraint (\ref{7/7/21.1}) of Definition \ref{5/22/25.3} will be the right-hand-side of (\ref{5/7/22.20}) with $\gamma_1=1$.  Setting $\frac{1}{\hat{t}_{1,n}}t = (\gamma_n,\dots,\gamma_2, 1)$, the corollary follows.
\end{proof}

\begin{corollary}\label{5/17/22.1} For $s= (x_1,\dots,x_n)\in \mathbb{T}_n$ with $x_1\ne 0$, the usual norm is $x_1^n$.  The special unital norm family of $\mathbb{T}_n$ is the $(n-1)$-parameter family,
\begin{equation} \label{5/5/22.6} \ell_{\gamma_2,\dots,\gamma_n}(x_1,\dots,x_n) = x_1 e^{P_{\gamma_2,\dots,\gamma_n}\left(\frac{x_2}{x_1},\dots,\frac{x_n}{x_1}\right)},\end{equation} where $P_{\gamma_2,\dots,\gamma_n}(w_1,\dots,w_{n-1})$ is a real $(n-1)$-degree multivariate polynomial with coefficients dependent on the $(n-1)$ real parameters $\gamma_2,\dots,\gamma_n$.    \end{corollary}

\begin{proof} The usual norm of $s=(x_1,\dots,x_n)\in \mathbb{T}_n$ is immediately seen to be $x_1^n$.

With respect to the notation $G_n,H_n$ introduced at the beginning of the proof of Theorem \ref{5/4/22.2}, each $(\gamma_n,\dots,\gamma_2,1)\in\mathbb{T}_n$ is associated with a Toeplitz matrix $G_n^{(1)}$ having its $(1,n)$-entry equal to $1$.  Corollary \ref{7/1/22.1} indicates that $H^{(1)}_n\equiv G_n^{(1)}\slashed{\bf 1}_n$ is a normalized uncurling metric. It can now be appreciated from (\ref{11/7/20.2}) that  the associated unital norm is $\text{exp}(\left[\int_{\bf 1}^s (H_n^{(1)} t^{-1})\cdot \mathbf{d}t\right])$.  Elementary integration and a simple induction argument quickly results in (\ref{5/5/22.6}) and the $(n-1)$-degree of the multivariate polynomial in question.
\end{proof}

\vspace{.75cm}

  \section{Extension to certain not associative algebras}
There are three isomorphism classes of two-dimensional real unital algebras, represented by $\mathbb{C}$, $\slashed{\mathbb{C}}$, and $\mathbb{D}$.  The complex numbers $\mathbb{C}$ can be thought of as being generalized by the Cayley-Dixon sequence of algebras.  The split-complex numbers $\slashed{\mathbb{C}}$ can be thought of as being generalized by the formally real Spin Factor Jordan algebras.  The dual numbers $\mathbb{D}$ can be thought of as being generalized by the real upper triangular matrix algebras with at least some redundant diagonal components.  These generalizations of $\mathbb{D}$ are associative algebras, and have been treated in Appendix A and some examples in Section 7.  However, the above generalizations of $\mathbb{C}$ and $\slashed{\mathbb{C}}$ are not associative (apart from the quaternions).  Application of our program to these latter algebras will be presented in this appendix. 

Spin Factor Jordan Algebras and Cayley-Dixon Algebras after the quaternions are power-associative but not associative.  While the usual norm is thereby inapplicable, there is an ``Algebraic Norm" replacement \cite{jacobson:1963} (based on the minimal polynomial, and thus analogous to the usual norm as given by the constant term of a characteristic polynomial).  For the Spin Factor Jordan Algebras and Cayley-Dixon Algebras, the Algebraic Norms are essentially the  Minkowski  norm and Euclidean norm, respectively.  Thus, useful norms are not restricted to associative (or alternative) algebras, even though it turns out that the multiplicative property (i.e., $N(s_1s_2)=N(s_1)N(s_2)$) must then be sacrificed  \cite{schafer:1963}. 
So, despite a strong bias against a norm not satisfying the multiplicative property, these examples indicate that ``algebraic" norms that are not multiplicative can nevertheless be very significant \cite{mccrimmon:2004}. 

Of course, apart from the omnipresent unital norm group member comparable to the usual norm as guaranteed by Theorem \ref{11/17/22.1} for real finite-dimensional unital associative algebras (as well as all unital norms resulting from two-dimensional real unital algebras), unital norms are not typically multiplicative.  The multiplicative property is instead replaced by (\ref{5/23/23.3}).  The latter equation also replaces the multiplicative property for the Algebraic Norm in the cases of the not associative Spin Factor Jordan Algebras and Cayley-Dixon Algebras. 
Like the Algebraic Norm, the unital norm/uncurling metric constructions remain valid for these latter two classes of algebras.  The way we deal with these particular not associative algebras is as follows.

\subsection{Unital norms for $^*$-algebras whose Hermitian elements are real} 

Consider an algebra associated with an involutive antiautomorphism ``
$^*$ " (known as a $^*$-algebra, star-algebra, or involutive algebra) such that the Hermitian elements are members of $\mathbb{R}$ (thus, $s+s^*$ and $s^*s$ are both real; also, $s^*s = ss^*$ because $(s+s^*)s = s(s+s^*)$ since $(s+s^*)$ is real). 
Then \[ X^2 - (s+s^*)X + s^*s, \] is the minimal polynomial of any not real element $s$, and annihilates all elements, real and not real.  Consistent with comments in the second paragraph of this appendix, we define its constant term to be the ``Algebraic Norm" of an element $s$, by analogy with the characteristic polynomial defined for an associative algebra, which supplies the norm from the constant term of that element-associated polynomial annihilating all elements.


On these $^*$-algebras, the multiplicative identity $\mathbf{1}$ is $1$ (where the latter is thought of as the algebra element that is the multiplicative identity on the subalgebra $\mathbb{R}$ rather than a field element).  Accordingly, for these algebras, we define the expression $\|{\bf 1}\|^2$ (e.g., that appears in Definition \ref{9/21/23.2}) to be $1$.  We also have,

\begin{lemma}\label{9/14/22.2} The inverse of a unit $s$ in a finite-dimensional $^*$-algebra whose Hermitian elements are real is unique, and given by $s^{-1}=\frac{s^*}{ss^*}$.\end{lemma}

\begin{proof} 
We can write $s=\frac{s+s^*}{2} +\frac{s-s^*}{2}$, where the first  term on the right-hand-side is real and the second term on the right-hand-side is antisymmetric with respect to application of the $^*$-operation.  We define $\sigma$ to be the first term on the right-hand-side, and $\bar{s}$ to be the second term on the right-hand-side.  Suppose $x = \chi + \bar{x}$ is an inverse of $s$, where $\chi$ is real, and $\bar{x}$ is antisymmetric. Then \begin{equation}\label{9/14/22.1} sx = \sigma\chi + \sigma\bar{x}+\chi\bar{s} + \bar{s}\bar{x} = 1,\end{equation} and $xs = \sigma\chi + \sigma\bar{x}+\chi\bar{s} + \bar{x}\bar{s} = 1$.  Thus, $\bar{s}\bar{x} = \bar{x}\bar{s}$.  Since $\bar{s}^* = (\frac{s-s^*}{2})^* = -(\frac{s-s^*}{2}) = -\bar{s}$, we have $(\bar{s}\bar{x})^* = \bar{x}^*\bar{s}^* = \bar{x}\bar{s} = \bar{s}\bar{x}$.  Hence, $\bar{s}\bar{x}$ is real.  It follows from (\ref{9/14/22.1}) that $\sigma\chi+\bar{s}\bar{x} =1$ and $\sigma\bar{x}+\chi\bar{s} = 0$.  The latter has has a unique solution, yielding $\chi+\bar{x} = \frac{\sigma-\bar{s}}{\sigma^2-\bar{s}\bar{s}} = \frac{s^*}{ss^*}$, where the second equality follows since the right-hand-side is obviously an inverse ($1=s\frac{s^*}{ss^*} = \frac{s^*}{ss^*}s$ because $s^*s=ss^*$ is real and thus self-adjoint), and the inverse has just now been seen to be unique.
\end{proof}

Extending Definition \ref{7/21/23.1} to include this class of $^*$-algebras, we can now provide a special unital norm  by identifying a linear transformation $L$ satisfying (\ref{5/22/25.3}), (\ref{7/7/21.1}) and then evaluating (\ref{11/7/20.2}).  
 
 \begin{proposition}\label{1/11/22.1} Given a $^*$-algebra with vector space of elements $\mathbb{R}^n$ and whose Hermitian elements are members of $\mathbb{R}$, the Algebraic Norm on the $^*$-algebra in a sufficiently small neighborhood of ${\bf 1}$ is the square of a special unital norm. \end{proposition}

\begin{proof} By definition, the Algebraic Norm is \begin{equation}\label{1/15/22.1} N(s)= ss^* = N(s^*)\in\mathbb{R}.\end{equation} 
 We then have $N(\alpha s) = \alpha^2 N(s)$ for $\alpha\in \mathbb{R}$, i.e., degree-2 positive homogeneity of $N(s)$.  
For $N(s)\ne 0$, Lemma \ref{9/14/22.2} implies 
\begin{equation}\label{1/12/22.1} s^{-1} = \frac{s^*}{N(s)}.\end{equation}
Applying $N(\cdot)$ to both sides of  (\ref{1/12/22.1}), the degree-2 positive homogeneity and (\ref{1/15/22.1}) imply $N(s)N(s^{-1})=1$, so that \begin{equation}\label{1/15/22.2} \sqrt{N(s)}\sqrt{N(s^{-1})}=1. \end{equation}

We are required to show that $\sqrt{N(s)}$ is a special unital norm.  Thus, we must first verify that $\sqrt{N(s)}$ satisfies the requirements of the extended  version of Definition \ref{5/22/25.3} as indicated prior to the statement of this proposition. 

To begin with, (\ref{1/15/22.2}) already satisfies  (\ref{5/23/23.3}).
 Equation (\ref{1/15/22.1}) indicates that $N(s)$ is a quadratic form, so there is an orthonormal basis with respect to which the members of the algebra can be written as $s= (s_1,\dots, s_i,\dots)$ with $N(s) = \sum_i \delta_i s_i^2$ and $\delta_i = -1$, $0$, or $1$, since $^*$ is an involution.  Let $K_1 \equiv {\rm diag}\{\delta_i\}$.   Then, \begin{equation}\label{2/6/22.1} N(s) = s\cdot(K_1 s).\end{equation}
If $s^{-1}$ exists (i.e., $N(s)\ne 0$), we then have \begin{equation}\label{2/6/22.2} \nabla\left[\sqrt{N(s)}\right] = \frac{K_1s}{\sqrt{N(s)}}.\end{equation}  Define $K_2$ to be the algebra's involutive antiautomorphism operation $^*$.  It follows that $K_1K_2K_2[\frac{s}{\sqrt{N(s)}}] = \nabla\left[\sqrt{N(s)}\right]$. Defining $L\equiv K_1K_2$, dividing both sides of the equation in the last sentence by $\sqrt{N(s)}$, and using (\ref{1/15/22.1}) and (\ref{1/12/22.1}), we then have \begin{equation}\label{1/12/22.2} Ls^{-1}  = \frac{\nabla\left[\sqrt{N(s)}\right]}{\sqrt{N(s)}} .\end{equation} This equation  is valid for $N(s)>0$, and the latter pertains in a sufficiently small neighborhood of ${\bf 1}$.  Thus,  (\ref{11/7/20.1}) is satisfied.  From (\ref{2/6/22.1}), (\ref{2/6/22.2}), and (\ref{1/12/22.2}), we have \[ s\cdot [Ls^{-1}] = s\cdot K_1\left[\frac{s}{N(s)}\right] = 1,\] which is (\ref{7/7/21.1}).  
 
 To establish that $\sqrt{N(s)}$ is a unital norm, we  now only need to verify that (\ref{5/22/25.3}) is satisfied, i.e., $d\left([Ls^{-1}]\cdot \mathbf{d}s\right)=0$.  According to (\ref{2/6/22.2}), (\ref{1/12/22.2}), and the definition  of $K_1$, we have $d\left([Ls^{-1}]\cdot \mathbf{d}s\right) = d\left(\sum_i \frac{\delta_is_i ds_i}{N(s)}\right)$.  Performing the requisite differentiation to obtain the component of the exterior derivative associated with a choice of $i$ and $j$, we find \[ \frac{\partial [Ls^{-1}]_j}{\partial s_i} - \frac{\partial [Ls^{-1}]_i}{\partial s_j} = -2\frac{(\delta_j s_j)(\delta_i s_i)}{N^2(s)} +
2\frac{(\delta_i s_i)(\delta_j s_j)}{N^2(s)} = 0,  \] where $\frac{\partial [Ls^{-1}]_j}{\partial s_i}$ is the partial derivative of the $j$-th component of $[Ls^{-1}]$ with respect to the $i$-th component of $s$, etc. Thus,  $d\left([Ls^{-1}]\cdot \mathbf{d}s\right) = 0$.

Consequently, $\sqrt{N(s)}$ satisfies the special unital norm requirements of the extended version of Definition \ref{5/22/25.3}.  
\end{proof}


\subsection{The Cayley-Dixon sequence}  From the proof of Proposition \ref{1/11/22.1}, it is easily shown that a normalized uncurling metric associated with an algebra of the Cayley-Dixon sequence (all of whose members are $^*$-algebras whose Hermitian elements are real) is given by ${\rm diag}\{1,-1,-1,\dots,-1\}$, with the associated special unital norm seen to be the Euclidean norm.  It is easily shown from results in \mbox{Section 7} that $\mathbb{C}$ is associated with a 1-parameter family of normalized uncurling metrics, which follows from the fact that the space of uncurling metrics is two dimensional.  But employing the $\frac{n(n-1)}{2}$ constraint equations with $n>2$, in the context of Lemma \ref{9/14/22.2}, a simple argument indicates that all other algebras of the sequence have only a 1-dimensional anti-rotor.  It follows that there is only the single normalized uncurling metric given above.  
 
\subsection{Formally real Spin Factor Jordan Algebras}  Consider the bilinear product ``$\,\bullet\,$" such that for
$\alpha,\beta\in \mathbb{R},\,\, {\sf a},{\sf b}\in \mathbb{R}^m$,  \begin{equation}  (\alpha+ {\sf a})\bullet (\beta + {\sf b}) \equiv
    [\alpha \beta + {\sf a} \cdot {\sf b}] + [ \beta {\sf a} + \alpha {\sf b}] \in \mathbb{R}\oplus \mathbb{R}^m. \label{12/24/18.1}\end{equation}  This makes the vector space $\mathbb{R}\oplus \mathbb{R}^m$ a unital algebra with the multiplicative identity ${\bf 1}$ as simply $1$ (where in this instance $1$ is understood to be an element of the algebra rather than a member of the field).  For $m\ge 2$, this is known as a (formally real) Spin Factor Jordan Algebra \cite{mccrimmon:2004}.


 There are two involutive antiautomorphisms on each of these commutative algebras, one of which is the identity and the other is the conjugate operation, $(\alpha+ {\sf a})^* = (\alpha- {\sf a})$. The latter is the only one of the two involutive antiautomorphisms  for which the Hermitian elements are real.  Applying Proposition \ref{1/11/22.1}, it is easily shown that the special unital norm is the Minkowski norm.  Similar to the case of the Cayley-Dixon algebras, it is easy to show that this is the only special unital norm.

\enlargethispage{1cm}
\section{Issues for future development} 
\begin{enumerate}
\item Settling the conjectures of Section 8 is of overriding interest.  If any of the conjectures are true, that result might be deep.  If any are false, it will be necessary to try to tame the resulting unruly/disorganized situation with some semblance of an enveloping theory.
\item Corollary \ref{8/18/23.2} is surely only a partial listing of the invariants accessible with this approach.  A full accounting of the invariants stemming from the anti-rotor and its dual is required.
\item Can this program be extended to real infinite-dimensional algebras?
\item The essential feature of this approach is its differentiable aspect.  Thus, it is reasonable to suppose that extension of our program to complex algebras might be fruitful. 
\item There are many potentially interesting combinatorial questions with a bearing on algebra structure that arise from the newly identified invariants in this paper (Section 5). \item What accessible structural information regarding an algebra does the image of the Skewer contain that is independent of the information obtainable from analysis of the kernel of the Skewer? 
\item This program can be applied to nonunital algebras where, of course, $f(s)=s^{-1}$ is excluded, and it is clearly not true that anti-rotors associated with distinct choices for $f(s)$ are in general the same.  Is the application to nonunital algebras worthwhile?
\end{enumerate}

\vspace{.75cm}



   \begin{thebibliography}{99}
   
\bibitem{eom:2024} Encyclopedia of Mathematics [Internet].  Ringel CM, Representation of an associative algebra; [cited June 16, 2025]; [about 4 screens]. Available from: \\ https://encyclopediaofmath.org/wiki/Representation\_of\_an\_associative\_algebra.
          
       \bibitem{belitskii:2005} Belitskii G, Lipyanski R, Sergeichuk VV.  Problems of classifying associative or Lie algebras and triples of symmetric or skew-symmetric matrices are wild.  Linear Algebra Appl.  2005;407:249-262.
       
           \bibitem{hochschild:1945} Hochschild G.  On the cohomology groups of an associative algebra.   Ann Math.  1945;46(1):58-67.
       
            \bibitem{gerstenhaber:1964} Gerstenhaber M.  On the deformation of rings and algebras.  Ann Math.  1964;79(1):59-103.
             
            \bibitem{loday:1998} Loday J. Cyclic Homology, Grundlehren der Mathematischen Wissenschaften, Volume 301.  New York (NY): Springer; 1998.              
            
            \bibitem{fgreensite:2025b} Greensite F (2025), A formulation of Euclidean geometry leading to recognition of a new class of algebra isomorphism invariants. arXiv:2209.14119

\bibitem{fgreensite:2025c} Greensite F (2022), Origin of the uncurling metric formalism from an Inverse Problems regularization method. arXiv:2209.14137      

\bibitem{bourbaki:1989} Bourbaki N (1989) Algebra I, Springer, Berlin, p. 543.   

\bibitem{birkhoff:1932} Birkhoff, GD.  A set of postulates for plane geometry (based on scale and protractors).  Ann Math.  1932;33:329–345.               


\bibitem{bannwarth:2015} Bannwarth I, El Din MS.  Probabilistic algorithm for computing the dimension of real algebraic sets. In: ISSAC'15 - 40th International Symposium on Symbolic and Algebraic Computation, July 2015; Bath UK. 2015 p. 37-44.  10.1145/2755996.2756670, hal-01152751.     
                 
\bibitem{kobayashi:2021} Kobayashi Y, Shirayanagi K, Tsukada M, Takahasi S.  A complete classification of three-dimensional algebras over $\mathbb{R}$ and $\mathbb{C}$.  AEJM.   2021;14(8):2150131. (25 pages).

              \bibitem{jacobson:1963} Jacobson N  (1963) Generic norm of an algebra. Osaka mathematics Journal, 15(1):25-50.
  
        \bibitem{schafer:1963}  Schafer R (1963) On forms of degree $n$ permitting composition. Journal of Mathematics and Mechanics, 12:777-792.       
      
       \bibitem{mccrimmon:2004} McCrimmon K (2004) A taste of Jordan algebras, Springer, New York.
 
      \end{thebibliography}
       
 \end{document} 
 
 
 
